% !TEX root = MosicaFE_Guidebook.tex
% ===========================================================================
%  MosicaFE_Guidebook.tex
%
%  MosicaFE user guide and developer manual, written after the template of
%  guidebook.pdf.
%
%  LANGUAGE POLICY (project convention, do not change):
%    * Every version of this guidebook is written in English.
%    * MosicaFE/README.md stays in Chinese.
%    * Code listings and comments inside this document are English too.
%
%  BUILD (no CJK text is present, so pdfLaTeX is enough; XeLaTeX also works):
%      pdflatex MosicaFE_Guidebook.tex
%      pdflatex MosicaFE_Guidebook.tex      % second pass: TOC and cross-refs
%
%  or:  latexmk -pdf MosicaFE_Guidebook.tex
% ===========================================================================
\documentclass[11pt,oneside]{book}

\usepackage[T1]{fontenc}
\usepackage[utf8]{inputenc}
\usepackage{lmodern}
\usepackage{amsmath}
\usepackage{amssymb}
\usepackage{amsthm}
\usepackage{booktabs}
\usepackage{longtable}
\usepackage{array}
\usepackage{graphicx}
\usepackage{listings}
\usepackage{upquote}
\usepackage{xcolor}
\usepackage{tcolorbox}
\usepackage{algorithm}
\usepackage{algorithmic}
\usepackage{enumitem}
\usepackage{caption}
\usepackage{geometry}
\usepackage{fancyhdr}
\usepackage[hidelinks]{hyperref}

\geometry{a4paper,left=2.6cm,right=2.6cm,top=2.8cm,bottom=2.8cm}

% ---------------------------------------------------------------------------
%  Code typesetting
% ---------------------------------------------------------------------------
\definecolor{codebg}{RGB}{248,248,248}
\definecolor{codegreen}{RGB}{0,110,60}
\definecolor{codegray}{RGB}{120,120,120}
\definecolor{codeblue}{RGB}{20,60,160}
\definecolor{codered}{RGB}{170,30,30}

\lstset{
  basicstyle=\ttfamily\small,
  keywordstyle=\color{codeblue}\bfseries,
  commentstyle=\color{codegreen}\itshape,
  stringstyle=\color{codered},
  numberstyle=\tiny\color{codegray},
  backgroundcolor=\color{codebg},
  frame=single,
  rulecolor=\color{codegray!40},
  framexleftmargin=4pt,
  numbers=left,
  numbersep=8pt,
  showstringspaces=false,
  breaklines=true,
  breakatwhitespace=false,
  tabsize=4,
  columns=fullflexible,
  keepspaces=true,
  captionpos=b,
  xleftmargin=18pt,
  framexrightmargin=0pt,
  upquote=true
}

\lstdefinelanguage{MosicaFE}{
  language=Python,
  morekeywords={as,with,None,True,False,lambda},
}

% Console output styles
\lstdefinestyle{shell}{
  language=,
  numbers=none,
  backgroundcolor=\color{black!4},
  basicstyle=\ttfamily\footnotesize
}
\lstdefinestyle{textout}{
  language=,
  numbers=none,
  backgroundcolor=\color{black!4},
  basicstyle=\ttfamily\footnotesize
}

% ---------------------------------------------------------------------------
%  Theorem-like environments
% ---------------------------------------------------------------------------
\theoremstyle{definition}
\newtheorem{definition}{Definition}[chapter]
\newtheorem{example}{Example}[chapter]
\newtheorem{remark}{Remark}[chapter]
\newtheorem{property}{Property}[chapter]
\theoremstyle{plain}
\newtheorem{theorem}{Theorem}[chapter]
\newtheorem{lemma}{Lemma}[chapter]

\newenvironment{important}{%
  \begin{tcolorbox}[colback=blue!5,colframe=blue!45,title=Key Point,fonttitle=\bfseries]%
}{%
  \end{tcolorbox}%
}

\newenvironment{caution}{%
  \begin{tcolorbox}[colback=orange!5,colframe=orange!55,title=Caution,fonttitle=\bfseries]%
}{%
  \end{tcolorbox}%
}

\newcommand{\mosica}{\textsf{MosicaFE}}
\newcommand{\code}[1]{\texttt{#1}}
\newcolumntype{L}[1]{>{\raggedright\arraybackslash}p{#1}}
\newcommand{\module}[1]{\texttt{MosicaFE.#1}}
\newcommand{\dif}{\,\mathrm{d}}

\captionsetup{font=small,labelfont=bf}
\renewcommand{\arraystretch}{1.12}

\title{\Huge \textbf{MosicaFE}\\[0.4em]
       \Large A Unified Python Library for FEM, VEM and RT0\\
       \Large Discretizations of Partial Differential Equations\\[0.6em]
       \large User Guide and Developer Manual}
\author{\normalsize MosicaFE contributors}
\date{\normalsize September 20, 2026}

% ===========================================================================
\begin{document}

\frontmatter
\maketitle

\begin{center}
\begin{minipage}{0.86\textwidth}
\small
\textbf{Abstract.} This manual documents \mosica{}, a Python library that brings the
finite element method (FEM), the virtual element method (VEM) and the lowest-order
Raviart--Thomas mixed element (RT0$\times$P0) under a single set of interfaces.
The user writes only a few lines in the driver file: choose a mesh, choose an
``element'', choose a linear solver, choose a visualization tool, and a
partial differential equation --- Poisson, reaction-diffusion, Helmholtz or the
RT0$\times$P0 mixed formulation --- is discretized and solved. The layered design
(\module{core}, \module{spaces}, \module{physics}, \module{solvers},
\module{postprocess}) makes both \textbf{adding a new element type} and
\textbf{adding a new equation} a local modification. The manual contains complete
usage examples, the mathematical background, verification of accuracy and
convergence, and templates for extending the framework.
\end{minipage}
\end{center}

\vspace{1.0em}

\tableofcontents

\mainmatter

% ===========================================================================
\chapter{Introduction}
\label{chap:intro}

\section{Overview}

\mosica{} is a partial differential equation solver aimed at teaching and research.
It merges three previously independent programs (classical finite elements, virtual
elements, and RT0 mixed finite elements) into \textbf{a single package with a single
set of interfaces}, sharing one mesh representation, one set of solvers and one set
of postprocessing tools.

The core design goals of the library are:

\begin{enumerate}[label=(\arabic*),itemsep=0.25em]
  \item \textbf{Minimal call syntax}: driver files are usually shorter than 10 lines;
  \item \textbf{Decoupled modules}: mesh, element, equation, linear algebra and
        postprocessing are mutually independent;
  \item \textbf{Easy extension}: a new element or a new PDE requires implementing
        only a small number of interfaces;
  \item \textbf{Verifiable}: convergence studies and error norms are built in, and
        the observed convergence orders can be reproduced with one command;
  \item \textbf{Self-contained}: the only hard dependency is \code{numpy}
        (\code{scipy} and \code{matplotlib} are optional).
\end{enumerate}

\section{Framework Philosophy}

The design of \mosica{} mirrors the mathematical structure of the finite element
method:

\begin{equation}
\underbrace{\text{PDE}}_{\text{physics}}
\;\xrightarrow{\ \text{weak form}\ }\;
\underbrace{\text{variational problem}}_{\text{spaces}}
\;\xrightarrow{\ \text{discretization}\ }\;
\underbrace{Au=b}_{\text{assembly}}
\;\xrightarrow{\ \text{solve}\ }\;
\underbrace{u_h}_{\text{postprocess}}
\label{eq:philosophy}
\end{equation}

Each step is the responsibility of a separate module, and modules communicate only
through explicit interfaces:

\begin{itemize}[itemsep=0.2em]
  \item \module{core} provides meshes and quadrature rules;
  \item \module{spaces} translates ``which element'' into a \textbf{matrix-level
        interface}: stiffness matrix, mass matrix, load vector (see
        Chapter~\ref{chap:spaces});
  \item \module{physics} only takes linear combinations of those matrices, so a new
        PDE usually takes a dozen lines;
  \item \module{solvers} is responsible for linear algebra;
  \item \module{postprocess} is responsible for errors, convergence and
        visualization.
\end{itemize}

\begin{important}
The orthogonal decomposition of ``element'' and ``equation'' is what distinguishes
\mosica{} most clearly from the three original programs. Those programs hard-wired
``triangle mesh $+$ linear element $+$ Poisson $+$ \code{spsolve} $+$ plotting''
into a single file; here these concerns are independent: changing the element does
not require touching the equation, and changing the equation does not require
touching the element.
\end{important}

\section{Quick Start}

Here is a complete 2D Poisson solver:

\begin{lstlisting}[language=MosicaFE,caption={Solving the Poisson equation in two lines},label={lst:quickstart}]
import numpy as np
from MosicaFE import Mesh, solve_poisson

source = lambda x: 2*np.pi**2 * np.sin(np.pi*x[0]) * np.sin(np.pi*x[1])

u = solve_poisson(Mesh.rectangle(nx=32, ny=32), element="P2", f=source)
print("degrees of freedom:", u.size)
\end{lstlisting}

The equivalent ``step-by-step'' form (which is also the standard form used
throughout this manual):

\begin{lstlisting}[language=MosicaFE,caption={Choosing mesh / element / solver / visualization},label={lst:quicksplit}]
import numpy as np
from MosicaFE import Mesh, FESpace, PoissonProblem
from MosicaFE.postprocess import plot_solution

def exact(x):
    return np.sin(np.pi*x[0]) * np.sin(np.pi*x[1])

def source(x):
    return 2*np.pi**2 * exact(x)

mesh = Mesh.rectangle(nx=32, ny=32)         # (1) choose a mesh
V = FESpace.lagrange(mesh, degree=2)        # (2) choose an element
problem = PoissonProblem(V, f=source, g=lambda x: 0.0)
u = problem.solve(method="direct")          # (3) choose a solver
plot_solution(u, mesh, V, filename="p2.png")  # (4) choose a visualization

print(problem.compute_errors(u, exact))
\end{lstlisting}

Replacing \code{FESpace.lagrange(mesh, degree=2)} in
Listing~\ref{lst:quicksplit} by \code{FESpace.vem(mesh)} no longer produces a finite
element solution: it produces the virtual element solution of the \textbf{same
equation} on the \textbf{same mesh}, and nothing else in the program changes.

\textbf{Changing the equation costs one line as well}. Built-in equations are the
Poisson equation, the reaction-diffusion equation, the Helmholtz equation and the
RT0$\times$P0 mixed formulation; each of them has a named constructor and a
one-line helper:

\begin{lstlisting}[language=MosicaFE,caption={Changing the equation},label={lst:quickswitch}]
from MosicaFE import Mesh, PoissonProblem, HelmholtzProblem, solve_helmholtz

mesh = Mesh.rectangle(nx=32, ny=32)

# -Laplacian(u) = f          (k = 0 limit of the line below)
u = PoissonProblem(V, f=source, g=zero).solve()

# -Laplacian(u) - k^2 u = f   (constant wavenumber k, Chapter 6.3)
u = HelmholtzProblem(V, f=source, g=zero, k=1.0).solve()
u = solve_helmholtz(mesh, element="P2", f=source, k=1.0)     # one line
\end{lstlisting}

\begin{important}
This is the reason the library exists: changing the discretization method costs one
line of code.
\end{important}

\section{Package Structure}
\label{sec:package-structure}

\begin{table}[htbp]
\centering
\caption{Module layout of \mosica{}}
\label{tab:modules}
\begin{tabular}{@{}lp{10.4cm}@{}}
\toprule
Module & Responsibility \\
\midrule
\module{core} & Mesh data structures and generation (interval / rectangle / box /
polygon / Voronoi), \textbf{facet topology and orientation} (\code{facets.py}, the
basis of $H(\mathrm{div})$-conforming spaces), quadrature rules, geometric tools
(area, volume, Jacobian, gradient transformation) \\
\module{spaces} & Discrete spaces: \code{LagrangeSpace} (P1/P2/Q1/Q2),
\code{VEMSpace} ($k=1$ virtual elements), \code{RT0Space} (RT0$\times$P0), and the
factory \code{FESpace} \\
\module{physics} & PDE weak forms and assembly: \code{PoissonProblem},
\code{MixedPoissonProblem}, \code{ReactionDiffusionProblem},
\code{HelmholtzProblem}, plus a PDE registry \\
\module{solvers} & Linear solvers: direct methods, CG, GMRES, BiCGSTAB, MINRES and
automatic selection \\
\module{postprocess} & Error norms ($L^2$ / $H^1$ / $L^\infty$), convergence studies,
visualization, and error norms for mixed formulations (pressure $+$ flux) \\
\bottomrule
\end{tabular}
\end{table}

\begin{important}
All three element families (FEM, VEM and RT0) are \textbf{implemented inside the
library itself}: the explicit RT0$\times$P0 basis functions, the exact mass matrix,
the facet topology and the saddle-point assembly live in \code{spaces/rt0.py},
\code{core/facets.py} and \code{physics/mixed\_poisson.py}, respectively.
\mosica{} neither depends on nor requires any external finite element backend.
\end{important}

\section{Organization of This Manual}

\begin{itemize}[itemsep=0.2em]
  \item Chapter~\ref{chap:install}: installation and environment;
  \item Chapter~\ref{chap:math}: mathematical background of the three methods;
  \item Chapters~\ref{chap:core}--\ref{chap:spaces}: meshes, quadrature and discrete
        spaces;
  \item Chapters~\ref{chap:physics}--\ref{chap:solvers}: equation assembly and
        linear solution;
  \item Chapters~\ref{chap:error}--\ref{chap:visual}: error analysis and
        visualization;
  \item Chapter~\ref{chap:examples}: complete examples;
  \item Chapters~\ref{chap:advanced}--\ref{chap:extend}: advanced topics and
        framework extension;
  \item Chapter~\ref{chap:faq}: frequently asked questions;
  \item Appendices~\ref{app:api}--\ref{app:bib}: API reference, mathematical
        notation, bibliography.
\end{itemize}

% ===========================================================================
\chapter{Installation and Setup}
\label{chap:install}

\section{System Requirements}

\subsection{Python Version}

Python $\ge$ 3.9 is required. The environment used for development and verification is
Python 3.12.10 (conda environment \code{fealpy}), with \code{numpy} 2.2.4,
\code{scipy} 1.15.2 and \code{matplotlib} 3.10.1.

\subsection{Dependencies}

\begin{table}[htbp]
\centering
\caption{Dependencies}
\label{tab:deps}
\begin{tabular}{@{}lll@{}}
\toprule
Package & Required & Purpose \\
\midrule
\code{numpy} & yes & arrays, dense linear algebra \\
\code{scipy} & optional & sparse matrices, \code{spsolve}, Krylov solvers, Voronoi meshes \\
\code{matplotlib} & optional & visualization \\
\bottomrule
\end{tabular}
\end{table}

If \code{scipy} is missing, \mosica{} automatically falls back to a purely
\code{numpy}-based dense implementation (handled uniformly in \module{\_compat}),
and small problems remain solvable. If \code{matplotlib} is missing, all plotting
functions print a notice and return \code{None} instead of raising an exception.

\section{Installation}

\subsection{Method 1: Installing as a Package (Recommended)}

\begin{lstlisting}[style=shell]
cd D:\mosica\MosicaFE
python -m pip install -e .
\end{lstlisting}

\subsection{Method 2: Using the Sources Directly}

It is enough to put the \textbf{parent directory} of \mosica{} on the module search
path; no installation is needed:

\begin{lstlisting}[style=shell]
# PowerShell
cd D:\mosica
$env:PYTHONPATH = "D:\mosica"
& "D:\anaconda\envs\fealpy\python.exe" MosicaFE\examples\poisson_demo.py
\end{lstlisting}

\section{Verifying the Installation}

\mosica{} ships two scripts that do not depend on a third-party test framework:

\begin{lstlisting}[style=shell]
& "D:\anaconda\envs\fealpy\python.exe" MosicaFE\tests\run_tests.py
& "D:\anaconda\envs\fealpy\python.exe" MosicaFE\tests\run_checks.py
\end{lstlisting}

The first runs 106 unit and integration tests; the second prints the measured
convergence orders of the various methods. An excerpt of the expected output:

\begin{lstlisting}[style=textout]
=== P2 ===
  level        h                  L2    rate          H1    rate        Linf    rate
      0  0.35355    3.4643e-02       -  4.4391e-01       -  1.4645e-01       -
      1  0.17678    4.4076e-03    2.97  1.1632e-01    1.93  3.8060e-02    1.94
      2  0.08839    5.5040e-04    3.00  2.9416e-02    1.98  9.6074e-03    1.99
      3  0.04419    6.8699e-05    3.00  7.3752e-03    2.00  2.4076e-03    2.00
  rate[L2] (last 3) = 3.002
  rate[H1] (last 3) = 1.990
\end{lstlisting}

If \code{pytest} is installed, one may also run

\begin{lstlisting}[style=shell]
python -m pytest MosicaFE/tests -q
\end{lstlisting}

\section{Compiling This Manual}

The document uses the standard \code{book} class and contains no CJK text, so
pdfLaTeX{} is sufficient; compile twice to obtain the table of contents and the
cross-references:

\begin{lstlisting}[style=shell]
pdflatex MosicaFE_Guidebook.tex
pdflatex MosicaFE_Guidebook.tex
\end{lstlisting}

Alternatively use \code{latexmk -pdf MosicaFE\_Guidebook.tex}. XeLaTeX{} produces the
same output. TeXShop users select \code{pdfLaTeX} as the engine and press Typeset.

% ===========================================================================
\chapter{Mathematical Background}
\label{chap:math}

\section{Model Problem}

The model problem solved by the library (and by the three original programs) is the
Poisson equation

\begin{equation}
\begin{cases}
-\nabla\cdot(\kappa\nabla u) = f, & \text{in } \Omega,\\[2pt]
u = g, & \text{on } \partial\Omega,
\end{cases}
\label{eq:poisson}
\end{equation}

where $\Omega\subset\mathbb{R}^d$ ($d=1,2,3$) is a bounded domain and $\kappa>0$ is
the diffusion coefficient. For $g=0$ and $\kappa=1$ this is the classical problem
$-\Delta u=f$.

\textbf{The Helmholtz model problem.}\label{sec:helmholtz-model}
The second model problem is the (frequency-domain, time-harmonic) \textbf{Helmholtz
equation} with a \emph{constant} wavenumber $k$ and a constant coefficient
$\kappa$,

\begin{equation}
\begin{cases}
-\nabla\cdot(\kappa\nabla u)-k^2u=f, & \text{in } \Omega,\\[2pt]
u=g, & \text{on } \partial\Omega .
\end{cases}
\label{eq:helmholtz}
\end{equation}

It is the equation solved by the built-in \code{HelmholtzProblem} and by the
one-line helper \code{solve\_helmholtz}. Two remarks are needed.

\begin{important}
\textbf{Low frequency means non-resonant.} The operator is
$-\Delta-k^2$. On a bounded domain with Dirichlet data, $-\Delta$ has a discrete
spectrum $0<\lambda_1\le\lambda_2\le\dots$; the solution of \eqref{eq:helmholtz} is
unique and the bilinear form
$a(u,v)=\kappa(\nabla u,\nabla v)-k^2(u,v)$ is positive definite --- and so is the
discrete matrix $A=\kappa K-k^2M$ --- if and only if $k^2<\lambda_1$. On the unit
square $\lambda_1=2\pi^2\approx19.74$, so the low-frequency range is
$k<\sqrt2\,\pi\approx4.443$. As soon as $k^2$ crosses $\lambda_1$ (or any other
eigenvalue) the matrix becomes \textbf{symmetric indefinite} and ill-conditioned:
the linear system must be solved by \code{"direct"} or \code{"minres"}, and the
error is amplified. A discrete, mesh-dependent version of the criterion is
$\lambda_{1,h}$, the smallest eigenvalue of the generalized problem
$Kv=\lambda Mv$ on the interior degrees of freedom; by the Galerkin theory
$\lambda_{1,h}\ge\lambda_1$ and $\lambda_{1,h}\to\lambda_1$ as $h\to0$.
\end{important}

\begin{caution}
\textbf{A second condition has no analogue for the Poisson problem.} The discrete
wave must be resolved by the mesh: the phase error (numerical dispersion) of the
scheme is controlled by $kh$, and for $kh\gtrsim1$ the error is dominated by
\emph{pollution} rather than by interpolation. Refining the mesh only restores the
optimal order once $kh\lesssim1$; a convergence table for the Helmholtz problem
must therefore always be read together with the $kh$ column (Chapter~\ref{chap:physics}
assembles the equation, Chapter~\ref{chap:examples} reports the measured orders).
\end{caution}

The wavenumber is written $k$ (lower case) throughout this manual to avoid a clash
with the symbol $K$ for a mesh cell (Appendix~\ref{app:notation}).

\section{Weak Formulation}

\begin{definition}[Weak formulation]
Let $V=\{v\in H^1(\Omega): v|_{\partial\Omega}=0\}$ and
$V_g=\{v\in H^1(\Omega): v|_{\partial\Omega}=g\}$. Find $u\in V_g$ such that
\begin{equation}
a(u,v)=\ell(v),\qquad \forall v\in V,
\label{eq:weak}
\end{equation}
where
\begin{equation}
a(u,v)=\int_\Omega \kappa\nabla u\cdot\nabla v\dif x,\qquad
\ell(v)=\int_\Omega f v\dif x .
\label{eq:bilinear}
\end{equation}
\end{definition}

The right-hand side $\ell$ may also carry a Neumann data term
$\int_{\Gamma_N} h v\dif s$; the physics layer of \mosica{} reserves an extension
point for it.

\section{Finite Element Discretization (FEM)}

Let $V_h\subset V$ be a finite-dimensional subspace with basis
$\{\varphi_i\}_{i=1}^n$. The discrete problem reads: find
$u_h=\sum_j u_j\varphi_j$ such that

\begin{equation}
\sum_{j=1}^n u_j\,a(\varphi_j,\varphi_i)=\ell(\varphi_i),\qquad i=1,\dots,n,
\label{eq:fem-system}
\end{equation}

that is, the linear system
\begin{equation}
A u = b,\qquad
A_{ij}=\int_\Omega\kappa\nabla\varphi_j\cdot\nabla\varphi_i\dif x,\qquad
b_i=\int_\Omega f\varphi_i\dif x .
\label{eq:Au}
\end{equation}

Computations on the reference element $\hat K$ make use of the gradient
transformation

\begin{equation}
\nabla_x\varphi_i = J_K^{-\top}\nabla_{\hat x}\hat\varphi_i,
\qquad
\int_K g\dif x=\int_{\hat K}g\circ F_K\,|\det J_K|\dif\hat x .
\label{eq:transform}
\end{equation}

\begin{table}[htbp]
\centering
\caption{Lagrange elements available in \mosica{}}
\label{tab:lagrange}
\begin{tabular}{@{}llcc@{}}
\toprule
Element & Name & Local DOFs per cell & DOF type \\
\midrule
triangle & P1 & 3 & vertices \\
triangle & P2 & 6 & vertices $+$ edge midpoints \\
quadrilateral & Q1 & 4 & vertices \\
quadrilateral & Q2 & 9 & vertices $+$ edge midpoints $+$ cell centre \\
tetrahedron & P1 & 4 & vertices \\
tetrahedron & P2 & 10 & vertices $+$ edge midpoints \\
hexahedron & Q1 & 8 & vertices \\
hexahedron & Q2 & 27 & vertices $+$ edges $+$ faces $+$ cell interior \\
\bottomrule
\end{tabular}
\end{table}

\section{Virtual Element Discretization (VEM)}
\label{sec:vem-math}

The key idea of the virtual element method is that \textbf{the basis functions need
not be written down explicitly}. The degree-one ($k=1$) virtual element space is
defined as

\begin{equation}
V_h(K)=\Bigl\{v\in H^1(K):\ \Delta v\in\mathcal{P}_1(K),\
v|_e\in\mathcal{P}_1(e)\ \forall e\subset\partial K\Bigr\},
\label{eq:vem-space}
\end{equation}

and its degrees of freedom are simply the function values at the vertices. The
bilinear form on a cell is split into two parts:

\begin{equation}
a_h(u_h,v_h)=\underbrace{a\bigl(\Pi^{\nabla}u_h,\Pi^{\nabla}v_h\bigr)}_{A_{\rm proj}}
+\underbrace{S_K\bigl((I-\Pi)u_h,(I-\Pi)v_h\bigr)}_{S_{\rm stab}} .
\label{eq:vem-split}
\end{equation}

\subsection{Projection Operator}

Take the scaled monomial basis $m_0=1$, $m_{k+1}=x_k-x_{K,k}$ (with $x_K$ the cell
centroid). The projection $\Pi^{\nabla}:V_h(K)\to\mathcal{P}_1(K)$ satisfies

\begin{equation}
\int_K\nabla\Pi^{\nabla}v\cdot\nabla q\dif x
=\int_K\nabla v\cdot\nabla q\dif x
=-\int_K v\,\Delta q\dif x+\int_{\partial K}v\,(\nabla q\cdot n)\dif s
=\int_{\partial K}v\,(\nabla q\cdot n)\dif s ,
\label{eq:vem-proj}
\end{equation}

where the last equality uses $\Delta q=0$ ($q$ being a linear polynomial). Since
$v|_e$ is linear on every edge, the boundary integral
$\int_e v\,(\nabla q\cdot n)\dif s$ is computed \textbf{exactly by the trapezoidal
rule}:

\begin{equation}
\int_e v\,(\nabla q\cdot n)\dif s
=\frac{|e|}{2}\Bigl(v(a)(\nabla q\cdot n)+v(b)(\nabla q\cdot n)\Bigr).
\label{eq:vem-edge}
\end{equation}

In matrix form, let $G=[\nabla m_0,\dots,\nabla m_d]\in\mathbb{R}^{d\times(d+1)}$,
$M=|K|\,G^{\top}G$ and
$R_{\alpha i}=\int_{\partial K}\varphi_i(\nabla m_\alpha\cdot n)\dif s$; then

\begin{equation}
\Pi=\text{pinv}(M)\,R\in\mathbb{R}^{(d+1)\times n}.
\label{eq:vem-pi}
\end{equation}

\begin{caution}
The normalization of the constant mode is essential. If one sets $\Pi_{0i}=1/n$
(the arithmetic mean of the vertex values), the projection reproduces linear
polynomials exactly only when the ``arithmetic mean of the vertices'' coincides with
the centroid --- true for triangles and squares, but false for general polygons. The
correct normalization follows from
$\ p(x_K)=\overline{d}-\nabla p\cdot(\overline{X}-x_K)$:

\begin{equation}
\Pi_{0\cdot}=\frac{1}{n}\mathbf{1}^{\top}
-(\overline{X}-x_K)\,\Pi_{1:d,\cdot},
\qquad \overline{X}=\frac{1}{n}\sum_i X_i .
\label{eq:vem-const}
\end{equation}

Ignoring this makes the stabilization penalize linear functions, which destroys
consistency: on irregular polygonal meshes the $L^2$ convergence order drops from 2
to below 0.5. See Section~\ref{sec:pitfall-vem} for details.
\end{caution}

\subsection{Stabilization}

By default \mosica{} uses the \textbf{symmetric positive semi-definite} stabilization

\begin{equation}
S_{\rm stab}=\sigma\,(I-\Pi_{\rm dof})^{\top}(I-\Pi_{\rm dof}),\qquad
\Pi_{\rm dof}=D\,\Pi,\quad
D_{i\cdot}=[1,\ X_i-x_K],
\label{eq:vem-stab}
\end{equation}

with the scaling factor

\begin{equation}
\sigma=\frac{\operatorname{tr}(A_{\rm proj})}{\operatorname{tr}\bigl((I-\Pi_{\rm dof})^{\top}(I-\Pi_{\rm dof})\bigr)}
\label{eq:vem-sigma}
\end{equation}

so that both terms have comparable magnitude (which also aligns their physical
dimensions). Passing \code{stabilization="legacy"} reproduces the skew form
$(I-\Pi_{\rm dof})/n$ of the early scripts (not recommended), while
\code{stabilization="none"} is intended for diagnostics only, since the pure
projection term is singular.

\subsection{Load Vector and Mass Matrix}

Virtual basis functions have no explicit expression, so $\int_K f\varphi_i$ is
approximated by \textbf{piecewise linear interpolation on a fan triangulation (a
tetrahedral subdivision in 3D)}. Splitting the cell from its centroid into
$T_1,\dots,T_n$, where $T_i$ has the three corners $c$ (the centroid), $V_i$ and
$V_{i+1}$ with barycentric coordinates $L_0,L_1,L_2$, gives

\begin{equation}
\int_K f\varphi_j\dif x
\approx c_0[j]\sum_i\int_{T_i} f L_0\dif x
+\int_{T_j}f L_1\dif x
+\int_{T_{j-1}}f L_2\dif x,
\label{eq:vem-load}
\end{equation}

where $c_0[j]$ is the interpolation of the basis function $\varphi_j$ at the
centroid (computed by linear least-squares reconstruction, which guarantees both
linear reproduction and the partition of unity). For triangular meshes this formula
is exact; for general polygons it retains the correct geometric weights.

\begin{caution}
The early scripts used $|K|/n$ for every vertex as an approximation of
$\int_K\varphi_i$. On quadrilateral and triangular cells that approximation happens
to be good enough, but on irregular polygons it strongly overestimates the vertices
on ``short edges'', which stalls convergence. This is precisely why the load vector
and the mass matrix were changed to formula~\eqref{eq:vem-load}.
\end{caution}

\section{RT0\texorpdfstring{$\times$}{x}P0 Mixed Discretization}

The mixed form writes the Poisson equation as a first-order system

\begin{equation}
q=-\nabla p,\qquad \nabla\cdot q=f\ \text{ in }\Omega,\qquad
p=0\ \text{ on }\partial\Omega .
\label{eq:mixed}
\end{equation}

Weak formulation: find $(q,p)\in RT0\times P0$ such that

\begin{equation}
\int_\Omega \nu\cdot q\dif x-\int_\Omega p\,\nabla\cdot\nu\dif x=0,
\qquad
\int_\Omega r\,\nabla\cdot q\dif x=\int_\Omega f r\dif x
\label{eq:mixed-weak}
\end{equation}

for all $(\nu,r)$. After discretization one obtains the saddle-point system

\begin{equation}
\begin{pmatrix} M & B^{\top}\\ B & 0\end{pmatrix}
\begin{pmatrix} \mathbf{Q}\\ \mathbf{P}\end{pmatrix}
=\begin{pmatrix} 0\\ \mathbf{f}\end{pmatrix},
\label{eq:saddle}
\end{equation}

which is \textbf{symmetric indefinite}; the conjugate gradient method therefore does
not apply.

\begin{table}[htbp]
\centering
\caption{Degrees of freedom and boundary conditions of RT0$\times$P0}
\label{tab:rt0}
\begin{tabular}{@{}lll@{}}
\toprule
Quantity & Location of DOFs & Count \\
\midrule
flux $q_h$ & one per cell facet (2D edge / 3D triangular face), shared by neighbouring cells & $\#$facets \\
pressure $p_h$ & one constant per cell & $\#$cells \\
\midrule
Dirichlet condition $p=0$ & \multicolumn{2}{l}{\textbf{natural} boundary condition of the mixed form, no elimination needed} \\
\bottomrule
\end{tabular}
\end{table}

\section{Error Estimates}

\begin{theorem}[C\'ea's lemma]
Let $u$ be the exact solution and $u_h$ the finite element solution. Then
$\|u-u_h\|_V\le C\inf_{v_h\in V_h}\|u-v_h\|_V$.
\end{theorem}

\begin{theorem}[A priori error estimate]
If $u\in H^{k+1}(\Omega)$ and polynomials of degree $k$ are used, then
\begin{equation}
\|u-u_h\|_{L^2}\le C h^{k+1}|u|_{H^{k+1}},\qquad
|u-u_h|_{H^1}\le C h^{k}|u|_{H^{k+1}} .
\label{eq:apriori}
\end{equation}
\end{theorem}

\begin{example}[Measured convergence orders of this library]
\begin{itemize}[itemsep=0.15em]
  \item P1 (triangles): $L^2$ is $O(h^2)$, $H^1$ is $O(h)$;
  \item P2 / Q2: $L^2$ is $O(h^3)$, $H^1$ is $O(h^2)$;
  \item VEM ($k=1$, arbitrary polygons): $L^2$ is $O(h^2)$, $H^1$ is $O(h)$;
  \item RT0$\times$P0: pressure $L^2$ is $O(h)$, flux $L^2$ is $O(h)$.
\end{itemize}
Chapter~\ref{chap:error} reports the measured data.
\end{example}

% ===========================================================================
\chapter{Core Module: Mesh and Quadrature}
\label{chap:core}

\section{Mesh Data Structure}

\subsection{Mesh Components}

A \code{Mesh} object stores the following information:

\begin{table}[htbp]
\centering
\caption{Main attributes of \code{Mesh}}
\label{tab:mesh-attrs}
\begin{tabular}{@{}L{2.8cm}L{1.9cm}L{9.9cm}@{}}
\toprule
Attribute & Type & Meaning \\
\midrule
\code{nodes} & \code{(n\_nodes, dim)} array & vertex coordinates \\
\code{elements} & \code{list[ndarray]} & vertex indices of each cell, \textbf{ragged} \\
\code{element\_type} & \code{str} & \code{interval}/\code{triangle}/\code{quad}/\code{tet}/\code{hex}/\code{polygon}/\code{polyhedron} \\
\code{boundary\_nodes} & \code{ndarray} & boundary vertex indices (determined automatically from the facet topology) \\
\code{dim} & \code{int} & spatial dimension \\
\bottomrule
\end{tabular}
\end{table}

\begin{important}
\code{elements} is deliberately a \textbf{ragged list}: triangles, quadrilaterals,
pentagons, Voronoi polygons, tetrahedra and hexahedra can therefore share a single
assembly loop. When all cells have the same number of vertices, a
\code{(n\_cells, k)} array view can be obtained through \code{mesh.uniform\_vertices}.
\end{important}

\subsection{Facets}

Codimension-one cell faces (facets) are the basis of boundary detection: in 1D a
facet is an endpoint, in 2D an edge, and in 3D a triangular or quadrilateral face.
\code{mesh.facets} returns all facets, \code{mesh.boundary\_facets} returns those
owned by exactly one cell, and their vertex set is where the Dirichlet boundary
lives.

\begin{property}
For a two-dimensional polygonal mesh, facets are exactly the edges of the polygons;
for hexahedral meshes they are expanded from a fixed table of six faces; for general
three-dimensional polyhedra a face table can be supplied explicitly through
\code{Mesh(..., faces=...)}.
\end{property}

\begin{caution}
The \code{mesh.facets} described above is a ``\textbf{vertex set used for boundary
detection}''; it is sufficient for Lagrange and virtual elements (whose degrees of
freedom live on the vertices). An $H(\mathrm{div})$-conforming mixed element (RT0),
however, additionally needs to know ``\textbf{which cells share a given face}'' and
``toward which cell the global normal of that face points''. The library therefore
provides a second layer of facet topology: \code{core/facets.py}. The next
subsection shows how to use it.
\end{caution}

\subsection{Facet Topology and Orientation: \code{core/facets.py}}
\label{sec:facets}

The RT0 degrees of freedom are \textbf{facet normal fluxes}
$\Phi_f=\int_f q\cdot n_f\dif s$, and two neighbouring cells must share the same
unknown. This requires an explicit global numbering and orientation of the facets.
The library makes this a first-class citizen of the \code{core} layer,
\code{FacetTopology} (\code{mesh.facet\_topology}, built lazily and cached):

\begin{table}[htbp]
\centering
\caption{Main attributes of \code{FacetTopology} (\code{t = mesh.facet\_topology})}
\label{tab:facettopo}
\begin{tabular}{@{}L{4.4cm}L{10.0cm}@{}}
\toprule
Attribute / method & Meaning \\
\midrule
\code{t.facets} & \code{(n\_facets, d)}: \textbf{sorted} vertex list of each facet (the global unique key) \\
\code{t.cell\_facets} & \code{(n\_cells, d+1)}: \code{[c,i]} is the global index of the $i$-th local facet of cell $c$ \\
\code{t.signs} & $\pm1$: whether the global normal points out of that cell \\
\code{t.normals} & unnormalized global normals (magnitude $=$ facet measure) \\
\code{t.counts} & number of adjacent cells: 2 for interior facets, 1 for boundary facets \\
\code{t.boundary\_facets()} & array of boundary facet indices (\code{t.n\_boundary} is the count) \\
\bottomrule
\end{tabular}
\end{table}

\begin{lstlisting}[language=MosicaFE,caption={Facet topology and orientation}]
from MosicaFE import Mesh

mesh = Mesh.box(nx=2, ny=2, nz=2, element_type="tet")
t = mesh.facet_topology

print(t.summary())          # cells=48, facets=120, interior=72, boundary=48
print(t.cell_facets[0])     # global indices of the 4 local facets of cell 0
print(t.signs[0])           # outward normal signs (+1 / -1)

# the two signs of an interior facet are necessarily opposite: this is
# H(div) conformity in its discrete form
c0, c1 = t.facet_cells[0]
print(t.signs[c0][list(t.cell_facets[c0]).index(0)] *
      t.signs[c1][list(t.cell_facets[c1]).index(0)])   # -1
\end{lstlisting}

\begin{important}
\code{build\_facet\_topology} rejects two kinds of input that would ``silently give
wrong results'': non-simplicial meshes (quadrilaterals / polygons / hexahedra), and
meshes with a facet shared by more than two cells (non-conforming partitions). This
is the technical reason why RT0 only accepts triangular and tetrahedral meshes.
\end{important}

\section{Mesh Generation}
\label{sec:mesh-generation}

All meshes are created by class methods of \code{Mesh}, and all of them use a
\textbf{counterclockwise (2D) / positive-volume (3D)} orientation.

\begin{lstlisting}[language=MosicaFE,caption={Generating the various mesh types},label={lst:meshgen}]
from MosicaFE import Mesh

m1 = Mesh.interval(a=0.0, b=1.0, nx=20)                # 1D interval

m2 = Mesh.rectangle(nx=8, ny=8)                        # triangles
m3 = Mesh.rectangle(nx=8, ny=8, element_type="quad")   # quadrilaterals

m4 = Mesh.box(nx=4, ny=4, nz=4)                        # tets (Kuhn split)
m5 = Mesh.box(nx=4, ny=4, nz=4, element_type="hex")    # hexahedra
m6 = Mesh.box(nx=3, ny=3, nz=3, pattern="center")      # 12 tets per cube

m7 = Mesh.pentagon(nx=8, ny=8)                         # pentagon mesh
m8 = Mesh.voronoi_polygon(n_points=100, seed=42)       # Voronoi polygons

for m in (m1, m2, m3, m4, m5, m7, m8):
    print(m.summary())
\end{lstlisting}

\begin{table}[htbp]
\centering
\caption{Mesh factory quick reference}
\label{tab:mesh-factories}
\begin{tabular}{@{}L{6.4cm}L{1.9cm}L{6.3cm}@{}}
\toprule
Factory & Domain & Cells \\
\midrule
\code{Mesh.interval(a,b,nx)} & $[a,b]$ & segments \\
\code{Mesh.rectangle(...,"triangle")} & rectangle & triangles \\
\code{Mesh.rectangle(...,"quad")} & rectangle & quadrilaterals \\
\code{Mesh.box(...,"tet")} & box & tetrahedra \\
\code{Mesh.box(...,"hex")} & box & hexahedra \\
\code{Mesh.pentagon(nx,ny)} & unit square & pentagons \\
\code{Mesh.voronoi\_polygon(...)} & unit square & Voronoi polygons \\
\code{Mesh.from\_arrays(...)} & user-defined & inferred automatically \\
\bottomrule
\end{tabular}
\end{table}

\begin{remark}[Kuhn subdivision]
By default \code{Mesh.box(...,element\_type="tet")} splits every small cube into 6
tetrahedra \textbf{sharing the main diagonal} (Kuhn subdivision). The triangulations
of the common face of two neighbouring cubes then agree, so the mesh is conforming.
\code{pattern="center"} additionally inserts the cube centre and splits into 12
tetrahedra --- the mesh used by the original 3D VEM scripts.
\end{remark}

\begin{remark}[Voronoi meshes]
When a Voronoi mesh is generated, mirror points are added in 8 directions outside
the domain so that the Voronoi cells at the boundary close, after which the
subdivision is clipped back to the domain. After clipping and sorting, coincident
neighbouring vertices are removed automatically, which avoids degenerate cells (a
polygon with repeated vertices would produce zero-area triangles in the fan
subdivision).
\end{remark}

\section{Mesh Properties and Geometric Quantities}

\begin{lstlisting}[language=MosicaFE,caption={Reading mesh information}]
mesh = Mesh.rectangle(nx=8, ny=8)

print(mesh.dim, mesh.n_nodes, mesh.n_cells)     # 2 81 128
print(mesh.get_mesh_size())                     # max cell diameter
print(mesh.cell_measures())                     # array of areas

v = mesh.get_cell_vertices(0)                   # (k, dim) coordinates
h = mesh.get_cell_diameter(0)
c = mesh.get_cell_centroid(0)
print(mesh.boundary_nodes[:5])
\end{lstlisting}

\section{Mesh Refinement}

\begin{lstlisting}[language=MosicaFE,caption={Uniform refinement}]
coarse = Mesh.rectangle(nx=4, ny=4)
fine = coarse.refine_uniform()
print(coarse.n_cells, fine.n_cells)     # 32 128
\end{lstlisting}

\code{refine\_uniform()} is available for structured meshes only (\code{interval} /
\code{rectangle} / \code{box} / \code{pentagon}): it doubles the number of cells in
every direction and regenerates the mesh, so the number of triangles is multiplied
by 4 and the number of hexahedra by 8. For unstructured meshes such as Voronoi
meshes it raises a \code{ValueError}.

\section{Reference Elements}

\begin{table}[htbp]
\centering
\caption{Reference element conventions}
\label{tab:refelem}
\begin{tabular}{@{}ll@{}}
\toprule
Cell type & Reference element \\
\midrule
\code{interval} & $[0,1]$ \\
\code{triangle} & $\{\xi\ge0,\eta\ge0,\xi+\eta\le1\}$ \\
\code{quad} & $[0,1]^2$, vertices ordered $(0,0),(1,0),(1,1),(0,1)$ \\
\code{tet} & standard tetrahedron, vertices $(0,0,0),(1,0,0),(0,1,0),(0,0,1)$ \\
\code{hex} & $[0,1]^3$: counterclockwise 4 vertices on $z=0$, then the matching 4 on $z=1$ \\
\bottomrule
\end{tabular}
\end{table}

Simplex cells use the affine map $x=v_0+J_K\hat x$; quadrilaterals and hexahedra use
bilinear / trilinear maps on $[0,1]^d$, whose Jacobian matrix depends on the
reference point $\hat x$ and must therefore be evaluated at every quadrature point.

\begin{lstlisting}[language=MosicaFE,caption={Geometric tools}]
import numpy as np
from MosicaFE.core.utils import (
    compute_jacobian, transform_gradient, apply_affine_map,
    polygon_area, polyhedron_volume,
)

quad = np.array([[0,0],[2,0],[2,2],[0,2]], float)
J, detJ = compute_jacobian(quad, np.array([0.5, 0.5]))
print(J, detJ)                      # [[2,0],[0,2]] 4.0
print(polygon_area(quad))           # 4.0
print(apply_affine_map([[0.5,0.5]], quad))   # [[1, 1]]
\end{lstlisting}

\begin{caution}
\textbf{The Jacobian convention.} \code{compute\_jacobian} returns the
\textbf{same convention} of Jacobian for both families of cells: columns are partial
derivatives of physical coordinates with respect to reference coordinates,

\begin{equation}
J_{ab}=\frac{\partial x_a}{\partial\xi_b},
\qquad
J=\big[\ \partial x/\partial\xi_1\ \ \cdots\ \ \partial x/\partial\xi_d\ \big],
\label{eq:jac-convention}
\end{equation}

and only with this convention are the gradient transformation
\code{transform\_gradient(dN, J)} (which uses $J^{-\top}$ internally) and the Newton
inverse map of quadrilaterals / hexahedra (which uses $J^{-1}$) correct.
\textbf{Do not write $dN^{\top}V$ when extending the code} --- that is the transposed
Jacobian. For axis-aligned rectangles and boxes $J$ is diagonal and symmetric, so
the error is completely hidden; but as soon as a mesh is \textbf{rotated or
distorted} ($J\ne J^{\top}$) the gradients, the stiffness matrix and the
interpolation are silently wrong. This library fixed exactly this issue on
2026-09-21 (bug number 7); the regression test is
\code{tests/test\_tensor\_geometry.py}, which checks $J$ against finite differences
and verifies that the errors are identical bit by bit under rotations of
$0^\circ/15^\circ/30^\circ/45^\circ$.
\end{caution}

\section{Quadrature Rules}

\begin{lstlisting}[language=MosicaFE,caption={Using quadrature rules}]
from MosicaFE.core.quadrature import QuadratureRule, cell_quadrature

quad = QuadratureRule.get_quadrature("triangle", order=3)
print(quad.n_points, quad.weights.sum())     # 4 0.5

f = lambda xi: xi[0]**2 + xi[1]**2
value = sum(quad.weights[q] * f(quad.points[q]) for q in range(quad.n_points))
print(value)                                  # 1/6

# map to a physical cell (fan subdivision for polygons / polyhedra is automatic)
mesh = Mesh.voronoi_polygon(n_points=20, seed=1)
pts, w = cell_quadrature(mesh, 0, order=3)
print(w.sum(), mesh.get_cell_measure(0))
\end{lstlisting}

\begin{table}[htbp]
\centering
\caption{Available quadrature rules (\code{order} is the degree of the polynomial to be integrated exactly)}
\label{tab:quad}
\begin{tabular}{@{}lll@{}}
\toprule
Cell & Rule & Remarks \\
\midrule
\code{interval} & Gauss--Legendre & $n=\lceil(order+1)/2\rceil$ points, normalized to $\int_0^1$ \\
\code{triangle} & barycentre / 3-point / 4-point / 6-point & exact for degree 1/2/3/4 respectively \\
\code{quad} & tensor-product Gauss & $n$ points per direction \\
\code{tet} & 1-point / 4-point / 5-point & exact for degree 1/2/3 respectively \\
\code{hex} & tensor-product Gauss & $n$ points per direction \\
\bottomrule
\end{tabular}
\end{table}

\begin{caution}
For tensor-product cells the quadrature order has to be understood
\textbf{per direction}. The integrand of the \code{Q1} stiffness matrix is quadratic
in each direction, so at least 2 Gauss points per direction are needed --- \mosica{}
automatically rounds the order up to \code{max(2, 2*p)}, which the user normally does
not have to specify. Earlier versions used one point too few here, which made the
computed convergence order of 3D hexahedral elements negative.
\end{caution}

% ===========================================================================
\chapter{Discrete Spaces}
\label{chap:spaces}

\section{Space Abstraction}

All spaces derive from \code{BaseFESpace} and expose a \textbf{matrix-level
interface} to the layers above:

\begin{table}[htbp]
\centering
\caption{Key methods of \code{BaseFESpace}}
\label{tab:space-api}
\begin{tabular}{@{}L{6.2cm}L{9.2cm}@{}}
\toprule
Method & Meaning \\
\midrule
\code{stiffness\_matrix()} & $A_{ij}=\int\nabla\varphi_j\cdot\nabla\varphi_i$ \\
\code{mass\_matrix()} & $M_{ij}=\int\varphi_j\varphi_i$ \\
\code{load\_vector(f)} & $b_i=\int f\varphi_i$ \\
\code{boundary\_dofs()} & indices of the Dirichlet degrees of freedom \\
\code{interpolate(f)} & interpolate a function onto the degrees of freedom \\
\code{evaluate(cell,u,pts)} & reconstruct a function from the degrees of freedom inside a cell and evaluate it \\
\code{evaluate\_gradient(cell,u,pts)} & as above, returning the gradient \\
\code{nodal\_values(u)} & map a degree-of-freedom vector back to the mesh nodes (for visualization) \\
\code{get\_cell\_dofs(i)} & global indices of the degrees of freedom of cell $i$ \\
\code{get\_dof\_coordinates(k)} & physical coordinates of degree of freedom $k$ \\
\code{eval\_basis(xi)} / \code{eval\_basis\_grad(xi)} & basis functions on the reference element (explicit-basis spaces only) \\
\bottomrule
\end{tabular}
\end{table}

\begin{important}
This set of interfaces is the reason why ``a new PDE takes a dozen lines'': the
physics layer only ever forms linear combinations of matrices and handles boundary
conditions; it never has to know about the cell loop.
\end{important}

\section{Lagrange Elements (Classical Finite Elements)}

\subsection{P1 Triangles}

On the reference triangle, with barycentric coordinates
$\lambda_0=1-\xi-\eta$, $\lambda_1=\xi$, $\lambda_2=\eta$:

\begin{equation}
\hat\varphi_0=1-\xi-\eta,\qquad
\hat\varphi_1=\xi,\qquad
\hat\varphi_2=\eta .
\label{eq:p1}
\end{equation}

\subsection{P2 Triangles}

\begin{align}
\hat\varphi_0&=\lambda_0(2\lambda_0-1),&
\hat\varphi_1&=\lambda_1(2\lambda_1-1),&
\hat\varphi_2&=\lambda_2(2\lambda_2-1),\notag\\
\hat\varphi_3&=4\lambda_0\lambda_1\ (\text{edge }0\text{-}1),&
\hat\varphi_4&=4\lambda_1\lambda_2\ (\text{edge }1\text{-}2),&
\hat\varphi_5&=4\lambda_2\lambda_0\ (\text{edge }2\text{-}0).
\label{eq:p2}
\end{align}

\subsection{Tensor-Product Elements}

Quadrilaterals and hexahedra use tensor products of the 1D Lagrange basis. For Q2,
for instance, with the 1D nodes $\{0,\tfrac12,1\}$:

\begin{equation}
N_0(t)=2t^2-3t+1,\qquad
N_1(t)=-4t^2+4t,\qquad
N_2(t)=2t^2-t,
\label{eq:q2-1d}
\end{equation}

the cell basis functions are $N_{ijk}(\xi,\eta,\zeta)=N_i(\xi)N_j(\eta)N_k(\zeta)$.
Since the nodes form a complete tensor grid, $\{N_{ijk}\}$ is automatically a
\textbf{nodal basis} dual to the nodal values.

\subsection{Global Numbering of the Degrees of Freedom}

\mosica{} uses a ``register'' that maps geometric entities to global degrees of
freedom:

\begin{itemize}[itemsep=0.15em]
  \item vertex degrees of freedom: keyed by the global vertex index;
  \item edge degrees of freedom: keyed by the \textbf{ordered pair} of endpoints;
  \item face degrees of freedom (3D): keyed by the ordered 4-tuple of face vertices;
  \item cell-interior degrees of freedom: keyed by the cell index.
\end{itemize}

Neighbouring cells therefore share edge / face degrees of freedom automatically, and
the global space is $H^1$-conforming.

\begin{lstlisting}[language=MosicaFE,caption={Creating spaces and inspecting the degrees of freedom}]
from MosicaFE import Mesh, FESpace

mesh = Mesh.rectangle(nx=8, ny=8)

V1 = FESpace.lagrange(mesh, degree=1)
V2 = FESpace.lagrange(mesh, degree=2)
print(V1.n_dofs, V2.n_dofs)      # 81 289

print(V2.get_cell_dofs(0))       # local -> global DOF map
print(V2.get_dof_coordinates(0))
print(V2.boundary_dofs()[:6])
\end{lstlisting}

\section{Virtual Element Spaces}

\begin{lstlisting}[language=MosicaFE,caption={Virtual element space (arbitrary polygons)}]
from MosicaFE import Mesh, FESpace

mesh = Mesh.voronoi_polygon(n_points=100, seed=0)
V = FESpace.vem(mesh)                       # k = 1, vertex DOFs only

print(V.n_dofs, mesh.n_nodes)               # equal
K = V.stiffness_matrix()                    # projection + stabilization
M = V.mass_matrix()
b = V.load_vector(lambda x: 1.0)

print("row sums of M equal b?",
      abs(b - M.sum(axis=1)).max() < 1e-12)
\end{lstlisting}

\begin{property}
The stiffness matrix of \code{VEMSpace} is
\begin{enumerate}[label=(\arabic*),itemsep=0.1em]
  \item symmetric (because the stabilization is written as a Gram matrix);
  \item positive semi-definite, with the constant functions in its kernel;
  \item linearly exact (applying the cell matrix to any vector of linear degrees of
        freedom gives the energy $\int_K|\nabla p|^2$).
\end{enumerate}
Item (3) is exactly the ``patch test'', whose importance is explained in
Section~\ref{sec:pitfall-vem}.
\end{property}

\begin{table}[htbp]
\centering
\caption{The \code{stabilization} options of \code{VEMSpace}}
\label{tab:stab}
\begin{tabular}{@{}lp{8.6cm}@{}}
\toprule
Value & Meaning \\
\midrule
\code{"trace"} (default) & $\sigma(I-\Pi_{\rm dof})^{\top}(I-\Pi_{\rm dof})$,
symmetric positive semi-definite, $\sigma$ given by the trace ratio. Recommended \\
\code{"legacy"} & $(I-\Pi_{\rm dof})/n$ from the early scripts, skew, diverges on
irregular polygons; use only to reproduce old results \\
\code{"none"} & the projection term $A_{\rm proj}$ alone. The cell matrices are
singular; for diagnostics only, never for solving \\
\bottomrule
\end{tabular}
\end{table}

\section{RT0\texorpdfstring{$\times$}{x}P0 Mixed Space}

\begin{lstlisting}[language=MosicaFE,caption={Mixed space}]
from MosicaFE import Mesh, FESpace

mesh = Mesh.rectangle(nx=16, ny=16)       # triangle mesh
V = FESpace.rt0(mesh)

print(V.n_flux, V.n_pressure, V.n_dofs)   # facets, cells, sum
print(V.cell_flux_dofs(0))                # local facet -> global flux DOF
print(V.facet_signs(0))                   # outward orientation signs
print(V.local_flux_mass(0))               # exact RT0 mass matrix of cell 0 (3x3)
\end{lstlisting}

\subsection{The Two Blocks Exposed by the Space Layer}

The interface of the RT0 space is \textbf{parallel} to that of \code{LagrangeSpace}
and \code{VEMSpace}; only the ``single elliptic operator'' is replaced by ``the two
blocks of a saddle-point problem'':

\begin{table}[htbp]
\centering
\caption{Matrix primitives of RT0$\times$P0 (combined by the physics layer)}
\label{tab:rt0-blocks}
\begin{tabular}{@{}L{5.4cm}L{2.9cm}L{6.1cm}@{}}
\toprule
Method & Shape & Meaning \\
\midrule
\code{V.flux\_mass\_matrix()} & $(\Phi,\Phi)$ & $M_{ij}=\int_K\varphi_i\cdot\varphi_j$, \textbf{closed form, exact}, no quadrature needed \\
\code{V.divergence\_matrix()} & $(P,\Phi)$ & $B_{K,i}=\int_K\nabla\cdot\varphi_i=\sigma_{K,i}$, always $\pm1$ \\
\code{V.pressure\_load\_vector(f)} & $(P,)$ & $\int_K f$ \\
\midrule
\code{V.interpolate\_pressure(f)} & $(P,)$ & cell average of $f$ ($L^2$ projection onto P0) \\
\code{V.interpolate\_flux(q)} & $(\Phi,)$ & $\Phi_f=\int_f q\cdot n_f\dif s$ \\
\code{V.evaluate\_flux(c,Q,pts)} & $(n_q,d)$ & reconstruct $q_h=\sum_i F_i\varphi_i$ in a cell \\
\code{V.locate\_cell(x)} & int & locate the cell containing a point via barycentric coordinates \\
\bottomrule
\end{tabular}
\end{table}

\begin{lstlisting}[language=MosicaFE,caption={Using the two blocks directly}]
from MosicaFE import Mesh, FESpace
from MosicaFE._compat import bmat

mesh = Mesh.rectangle(nx=16, ny=16)
V = FESpace.rt0(mesh)

M = V.flux_mass_matrix()          # (Phi,Phi)
B = V.divergence_matrix()         # (P,Phi)
f = V.pressure_load_vector(src)   # (P,)

A = bmat([[M, -B.T], [B, None]])  # [[M, -B^T], [B, 0]]
\end{lstlisting}

\begin{caution}
The RT0 space \textbf{does not provide} \code{stiffness\_matrix()} /
\code{mass\_matrix()} / \code{load\_vector()}, nor \code{interpolate()}: it
corresponds to a saddle-point problem and carries two fields at once, so there is no
``single operator'' and no ``single interpolation''. Calling those methods raises a
\code{NotImplementedError} whose message points to the method to use instead (for
example \code{mass\_matrix()} $\to$ \code{flux\_mass\_matrix()},
\code{load\_vector()} $\to$ \code{pressure\_load\_vector()}). In day-to-day use,
combine the space with \code{MixedPoissonProblem}.
\end{caution}

\section{The Space Factory \code{FESpace}}

\begin{lstlisting}[language=MosicaFE,caption={The three space families}]
from MosicaFE import Mesh, FESpace

mesh = Mesh.rectangle(nx=16, ny=16)

V_fem = FESpace.lagrange(mesh, degree=2)    # P2
V_vem = FESpace.vem(mesh, degree=1)         # VEM k=1
V_rt0 = FESpace.rt0(mesh)                   # RT0 x P0

# by name (used by solve_poisson(element=...))
print(FESpace.available())
V_any = FESpace.create(mesh, family="vem", degree=1)
\end{lstlisting}

\begin{table}[htbp]
\centering
\caption{String names of the ``elements'' (used by \code{solve\_poisson(element=...)})}
\label{tab:elem-names}
\begin{tabular}{@{}L{3.6cm}L{3.6cm}L{6.2cm}@{}}
\toprule
Name & Space family & Available meshes \\
\midrule
\code{"P1"} / \code{"P2"} & Lagrange & triangles, tetrahedra \\
\code{"Q1"} / \code{"Q2"} & Lagrange (tensor product) & quadrilaterals, hexahedra \\
\code{"VEM"} / \code{"VEM1"} & virtual elements, $k=1$ & arbitrary polygons / polyhedra \\
\code{"RT0"} / \code{"RT0-P0"} / \code{"MIXED"} & RT0$\times$P0 & triangles, tetrahedra \\
\bottomrule
\end{tabular}
\end{table}

% ===========================================================================
\chapter{Physics Problems and Assembly}
\label{chap:physics}

\section{The Physics Problem Interface}

All PDEs derive from \code{BasePhysics} and only have to implement one method:

\begin{lstlisting}[language=MosicaFE,caption={The \code{BasePhysics} interface},label={lst:basephysics}]
class BasePhysics(ABC):
    def __init__(self, space, f=None, g=None, **kwargs):
        self.space = space
        self.mesh = space.mesh
        self.n_dofs = space.n_dofs
        self.f = f or (lambda x: 0.0)      # source term
        self.g = g or (lambda x: 0.0)      # Dirichlet data
        self.boundary_dofs = space.boundary_dofs()

    @abstractmethod
    def assemble(self):
        """return (A, b) with boundary conditions applied"""

    def solve(self, method="auto", **kwargs):
        A, b = self.assemble()
        return Solver.solve(A, b, method=method, **kwargs)
\end{lstlisting}

\section{The Poisson Problem}

\begin{lstlisting}[language=MosicaFE,caption={Assembly of the Poisson problem (complete implementation)}]
@register_problem("poisson")
class PoissonProblem(BasePhysics):
    def __init__(self, space, f=None, g=None, kappa=1.0, **kwargs):
        super().__init__(space, f=f, g=g, **kwargs)
        self.kappa = kappa

    def assemble(self):
        A = self.space.stiffness_matrix()
        if float(self.kappa) != 1.0:
            A = float(self.kappa) * A
        b = self.space.load_vector(self.f)
        A, b = self.apply_boundary_conditions(A, b)
        return A, b
\end{lstlisting}

\begin{important}
The whole implementation is 8 lines long, and it holds \textbf{simultaneously for
Lagrange elements and for virtual elements}. This is the payoff of pushing the
``cell loop'' down into the \module{spaces} layer.
\end{important}

\section{The Helmholtz Problem}
\label{sec:helmholtz}

The Helmholtz problem \eqref{eq:helmholtz} differs from the Poisson problem by a
single term, and correspondingly its implementation differs by a single line: the
stiffness matrix is replaced by the \emph{indefinite} combination
$\kappa K-k^2M$.

\begin{lstlisting}[language=MosicaFE,caption={Assembly of the Helmholtz problem (complete implementation)},label={lst:helmholtz}]
@register_problem("helmholtz")
class HelmholtzProblem(BasePhysics):
    """-div(kappa grad u) - k**2 u = f"""

    def __init__(self, space, f=None, g=None, k=1.0, kappa=1.0, **kwargs):
        super().__init__(space, f=f, g=g, **kwargs)
        self.k = float(k)              # constant wavenumber
        self.kappa = float(kappa)      # constant diffusion coefficient

    def assemble(self):
        A = self.space.stiffness_matrix() - self.k**2 * self.space.mass_matrix()
        if self.kappa != 1.0:
            A = self.kappa * A
        b = self.space.load_vector(self.f)
        A, b = self.apply_boundary_conditions(A, b)
        return A, b
\end{lstlisting}

The class is a subclass of \code{BasePhysics}, so it inherits \code{solve},
\code{compute\_error}/\code{compute\_errors} and everything else unchanged; the
only thing that had to be supplied is the combination of the two matrices that the
space already computed. Consequently the Helmholtz problem accepts \textbf{every
H1-conforming space}, exactly as the Poisson problem does:

\begin{lstlisting}[language=MosicaFE,caption={The same equation on any element family}]
import numpy as np
from MosicaFE import FESpace, HelmholtzProblem, Mesh, solve_helmholtz

k = 1.0
exact = lambda x: np.sin(np.pi*x[0]) * np.sin(np.pi*x[1])
source = lambda x: (2*np.pi**2 - k**2) * exact(x)     # manufactured solution

mesh = Mesh.rectangle(nx=32, ny=32)

u = HelmholtzProblem(FESpace.lagrange(mesh, 2), f=source, g=lambda x: 0.0,
                     k=k).solve(method="direct")
u = HelmholtzProblem(FESpace.vem(mesh), f=source, g=lambda x: 0.0,
                     k=k).solve(method="direct")

u = solve_helmholtz(mesh, element="P2", f=source, k=k)   # one-line helper
u = run("helmholtz", mesh, element="VEM", f=source, k=k) # by registry name
\end{lstlisting}

\begin{important}
Three practical points, all of them consequences of \eqref{eq:helmholtz} rather
than of the library.

\begin{enumerate}[label=(\arabic*),itemsep=0.2em]
  \item \textbf{Choose the solver by the shift, not by the element.} For
        $k^2<\lambda_{1,h}$ the matrix $\kappa K-k^2M$ is symmetric positive
        definite, so \code{"auto"}, \code{"direct"} and \code{"cg"} are all usable.
        Above the first eigenvalue it is symmetric \emph{indefinite}: \code{"direct"}
        and \code{"minres"} still work, CG does not.
  \item \textbf{Report $kh$ next to the error.} The a priori estimate
        \eqref{eq:apriori} is an $h\to0$ statement; for a wave problem the
        pre-asymptotic (pollution) regime can be wide, so a convergence table is
        only meaningful together with the resolution indicator $kh$.
  \item \textbf{Constant coefficients only.} The library accepts a real constant
        $k$ and a real constant $\kappa$ and rejects anything else with an explicit
        error message rather than assembling a matrix silently. The variant
        $-\Delta u+k^2u=f$ (the screened Poisson equation) is the built-in
        \code{ReactionDiffusionProblem} with \code{reaction\_coeff=k**2}; complex
        wavenumbers (absorbing media) would require a complex-valued pipeline and
        are outside the current scope.
\end{enumerate}
\end{important}

Table~\ref{tab:helmholtz-verification} collects the measured verification data for
the constant-wavenumber case $k=1$; the discrete first eigenvalue was obtained as
the smallest eigenvalue of the generalized problem $Kv=\lambda Mv$ restricted to
the interior degrees of freedom.

\begin{table}[htbp]
\centering
\caption{Helmholtz verification ($k=1$, $u=\sin(\pi x)\sin(\pi y)$, constant wavenumber)}
\label{tab:helmholtz-verification}
\begin{tabular}{@{}lcccc@{}}
\toprule
Space / mesh & $L^2$ order & $H^1$ order & $k^2/\lambda_{1,h}$ & $kh$ range \\
\midrule
P1 / triangles    & 1.98 & 0.99 & 0.0505 & $0.25\to0.031$ \\
P2 / triangles    & 3.01 & 1.99 & 0.0507 & $0.25\to0.031$ \\
Q1 / quadrilaterals & 2.00 & 1.00 & 0.0507 & $0.25\to0.031$ \\
VEM $k=1$ / triangles & 1.98 & 0.99 & 0.0505 & $0.25\to0.031$ \\
\bottomrule
\end{tabular}
\end{table}

The measured orders coincide with the Poisson ones, which is the point: in the
low-frequency regime the Helmholtz problem is \emph{not} harder than the Poisson
problem, and the same elements, solvers and postprocessing apply. The discrete
eigenvalues $\lambda_{1,h}=19.7868$ (P1, VEM) and $19.7394$ (P2) bracket the exact
value $2\pi^2=19.7392$ from above, as the Galerkin eigenvalue theory predicts.

\section{Boundary Conditions}

\subsection{Imposing Dirichlet Conditions Correctly}

To impose $u_j=g_j$ on the boundary degrees of freedom $j$, the standard approach is
``eliminate the row and the column, and correct the right-hand side'':

\begin{enumerate}[label=(\arabic*),itemsep=0.15em]
  \item move the contribution of the boundary columns to the right-hand side:
        \begin{equation}
        b_i \leftarrow b_i-\sum_{j\in\mathcal{D}}A_{ij}g_j,
        \qquad \forall i\notin\mathcal{D};
        \label{eq:bc-lift}
        \end{equation}
  \item delete all rows and columns involving $\mathcal{D}$;
  \item set $A_{jj}=1$ and $b_j=g_j$ for $j\in\mathcal{D}$.
\end{enumerate}

\begin{caution}
Step (1) ``looks'' dispensable for homogeneous conditions $g=0$, which is why many
demo programs leave it out; as soon as the boundary data are nonzero the solution is
wrong. \mosica{} guards against a regression of this kind with a simple but sensitive
test --- the \textbf{inhomogeneous linear patch test}:

\begin{equation}
u(x)=1+2x_1-\tfrac12 x_2,\qquad f\equiv0,\qquad g=u|_{\partial\Omega},
\label{eq:patch}
\end{equation}

where the exact solution belongs to the discrete space, so the numerical solution
must be \textbf{exact} (error $\sim10^{-15}$). Before the fix, the measured error was
$2.875$.
\end{caution}

\subsection{Other Boundary Conditions}

\begin{itemize}[itemsep=0.2em]
  \item \textbf{Neumann conditions}: add $\int_{\Gamma_N}h\varphi_i\dif s$ to the
        result of \code{load\_vector}; this requires a quadrature rule on the
        boundary facets;
  \item \textbf{RT0 mixed form:} $p=0$ is a natural boundary condition, so no degree
        of freedom has to be eliminated (which is also why mixed formulations are
        more ``natural'' for flux problems);
  \item inhomogeneous data $p=g\ne0$ in the mixed form require an additional boundary
        flux term; the current version reports this explicitly instead of returning a
        wrong result.
\end{itemize}

\section{Assembly Algorithm}

\begin{algorithm}[htbp]
\caption{Global assembly (\code{BaseFESpace.\_assemble} in \mosica{})}
\label{alg:assembly}
\begin{algorithmic}[1]
\REQUIRE space $V$, cell-local matrix generator \code{local\_matrix(cell)}
\STATE initialize the triplet lists \code{rows, cols, vals}
\FOR{each cell $c=0,\dots,n_{\rm cells}-1$}
    \STATE $\mathcal{D}_c\leftarrow$\code{get\_cell\_dofs(c)}
    \STATE $K_c\leftarrow$\code{local\_matrix(c)}
    \STATE append the flattened row / column indices of $\mathcal{D}_c$ and $K_c$ to
           the triplet lists
\ENDFOR
\STATE build the matrix from the COO triplets (repeated indices are summed)
\IF{symmetry is required}
    \STATE $A\leftarrow\frac12(A+A^{\top})$
\ENDIF
\RETURN $A$
\end{algorithmic}
\end{algorithm}

\begin{caution}
The right-hand side must be assembled with ``per-entry accumulation'' semantics
(\code{numpy}'s \code{np.add.at}). If one writes \code{b[dofs] += local} directly,
the buffered assignment of \code{numpy} silently drops accumulations whenever
\code{dofs} contains repeated indices (for instance when the mesh has duplicate
vertices), while the COO assembly of the sparse matrix sums them correctly; this
inconsistency produces errors that are very hard to notice.
\end{caution}

\section{Using the Poisson Problem}

\begin{lstlisting}[language=MosicaFE,caption={Assembling and solving}]
import numpy as np
from MosicaFE import Mesh, FESpace, PoissonProblem

def source(x):
    return 2*np.pi**2 * np.sin(np.pi*x[0]) * np.sin(np.pi*x[1])

mesh = Mesh.rectangle(nx=32, ny=32)
V = FESpace.lagrange(mesh, degree=2)

problem = PoissonProblem(V, f=source, g=lambda x: 0.0, kappa=1.0)

A, b = problem.assemble()
print("system size:", A.shape, "non-zeros:", A.nnz)

u = problem.solve(method="direct")
print("solution size:", u.size)
\end{lstlisting}

\section{The PDE Registry}

\begin{lstlisting}[language=MosicaFE,caption={Creating and solving equations by name}]
from MosicaFE import FESpace, Mesh, available_problems, create_problem, run

print(available_problems())        # ['mixed_poisson', 'poisson', ...]

mesh = Mesh.rectangle(nx=32, ny=32)
problem = create_problem("poisson", FESpace.vem(mesh),
                         f=source, g=lambda x: 0.0)
u = problem.solve()

u2 = run("poisson", mesh, element="VEM", f=source, g=lambda x: 0.0)
\end{lstlisting}

% ===========================================================================
\chapter{Linear Solvers}
\label{chap:solvers}

\section{The Solver Interface}

\begin{lstlisting}[language=MosicaFE,caption={The \code{Solver} interface}]
import numpy as np, time
from MosicaFE import FESpace, Mesh, PoissonProblem, Solver

mesh = Mesh.rectangle(nx=32, ny=32)
problem = PoissonProblem(FESpace.lagrange(mesh, 2), f=source,
                         g=lambda x: 0.0)
A, b = problem.assemble()

t0 = time.time(); u_direct = Solver.direct(A, b); t1 = time.time()
u_cg    = Solver.conjugate_gradient(A, b, tol=1e-10)
u_gmres = Solver.gmres(A, b, tol=1e-10, restart=30)
u_bicg  = Solver.bicgstab(A, b, tol=1e-10)

print("direct %.3f s, diff vs CG: %.3e"
      % (t1 - t0, np.linalg.norm(u_direct - u_cg)))
\end{lstlisting}

The problem object can also choose for itself:

\begin{lstlisting}[language=MosicaFE]
u = problem.solve(method="auto")   # auto/direct/cg/gmres/bicgstab/minres
\end{lstlisting}

\section{Direct Methods}

\begin{itemize}[itemsep=0.2em]
  \item sparse matrices: \code{scipy.sparse.linalg.spsolve} (sparse LU / Cholesky
        internally);
  \item dense matrices: \code{numpy.linalg.solve};
  \item complexity: about $O(n^{1.5})$ in 2D and $O(n^{2})$ in 3D.
\end{itemize}

\section{Iterative Methods}

\begin{table}[htbp]
\centering
\caption{Available Krylov solvers}
\label{tab:krylov}
\begin{tabular}{@{}llp{6.4cm}@{}}
\toprule
Method & Applicable matrices & Remarks \\
\midrule
\code{conjugate\_gradient} & symmetric positive definite & $O(\mathrm{nnz})$ per step, iteration count $O(\sqrt{\kappa})$ \\
\code{gmres} & general (nonsymmetric) & with restarts, storage $O(\text{restart}\times n)$ \\
\code{bicgstab} & general & two matrix-vector products per step, low memory \\
\code{minres} & symmetric indefinite & e.g.\ the RT0 saddle-point system \\
\bottomrule
\end{tabular}
\end{table}

\begin{caution}
The conjugate gradient method requires a symmetric positive definite matrix. The
stiffness matrix of virtual elements is symmetric under the default stabilization,
but it is no longer symmetric if \code{stabilization="legacy"} is selected, in which
case CG fails (possibly even without an error message, returning a wrong result).
For saddle-point systems use a direct method or MINRES.
\end{caution}

\section{Solver Selection Guide}

\begin{table}[htbp]
\centering
\caption{Solver selection recommendations}
\label{tab:solver-guide}
\begin{tabular}{@{}llll@{}}
\toprule
Problem size & Matrix type & Recommended solver & Remarks \\
\midrule
$<2\times10^4$ DOFs & any & direct & exact, no tuning required \\
$>2\times10^4$ & symmetric positive definite & CG & preconditioning recommended \\
$>2\times10^4$ & symmetric indefinite & MINRES or direct & saddle-point problems, and Helmholtz with $k^2>\lambda_1$ \\
$>2\times10^4$ & nonsymmetric & GMRES / BiCGSTAB & needs a restart parameter \\
\bottomrule
\end{tabular}
\end{table}

\begin{lstlisting}[language=MosicaFE,caption={Automatic selection}]
from MosicaFE import Solver

print(Solver.recommend(1000, symmetric=True))     # 'direct'
print(Solver.recommend(10**6, symmetric=True))    # 'cg'
print(Solver.recommend(10**6, symmetric=False))   # 'gmres'
\end{lstlisting}

% ===========================================================================
\chapter{Error Analysis and Convergence}
\label{chap:error}

\section{Error Norms}

\begin{align}
\|e\|_{L^2}&=\Bigl(\int_\Omega|u-u_h|^2\dif x\Bigr)^{1/2},&
\|e\|_{L^2}^2&\approx\sum_K\sum_q w_q\,|u(x_q)-u_h(x_q)|^2|\det J_K|,
\label{eq:l2norm}\\
|e|_{H^1}&=\Bigl(\int_\Omega|\nabla u-\nabla u_h|^2\dif x\Bigr)^{1/2},&
\|e\|_{L^\infty}&=\max_{x\in\Omega}|u(x)-u_h(x)| .
\label{eq:h1norm}
\end{align}

\begin{lstlisting}[language=MosicaFE,caption={Computing the three norms}]
import numpy as np

def exact(x):
    return np.sin(np.pi*x[0]) * np.sin(np.pi*x[1])

def exact_grad(x):
    return np.array([np.pi*np.cos(np.pi*x[0])*np.sin(np.pi*x[1]),
                     np.pi*np.sin(np.pi*x[0])*np.cos(np.pi*x[1])])

u = problem.solve()
l2   = problem.compute_error(u, exact, "L2")
linf = problem.compute_error(u, exact, "Linf")
h1   = problem.compute_error(u, exact, "H1", exact_grad=exact_grad)

print("%.6e %.6e %.6e" % (l2, h1, linf))

norms = problem.compute_errors(u, exact, exact_grad)
\end{lstlisting}

\begin{remark}
Virtual element basis functions have no explicit expression, so
\code{VEMSpace.evaluate} uses a ``linear least-squares reconstruction from the
vertex degrees of freedom''. The error of that reconstruction is of the same order
as $u_h$ itself, so the measured convergence order coincides with the true
convergence order of $u_h$ (see Section~\ref{sec:verify}).
\end{remark}

\section{Convergence Study Framework}

\begin{lstlisting}[language=MosicaFE,caption={Automatic convergence study},label={lst:convstudy}]
from MosicaFE import FESpace, Mesh, PoissonProblem, convergence_study
from MosicaFE.postprocess import estimate_convergence_rate

def mesh_factory(level):
    n = 4 * 2**level
    return Mesh.rectangle(nx=n, ny=n)

def problem_factory(mesh):
    V = FESpace.lagrange(mesh, degree=2)
    return PoissonProblem(V, f=source, g=lambda x: 0.0)

h, errors = convergence_study(mesh_factory, problem_factory,
                              exact, exact_grad,
                              refinement_levels=5, plot=True,
                              filename="conv_p2.png")

print(estimate_convergence_rate(h, errors["L2"], use_last=3))   # ~3.0
\end{lstlisting}

\section{Convergence Rate Estimation}

From the errors on two successive levels,

\begin{equation}
\text{rate}=\frac{\log(e_1/e_2)}{\log(h_1/h_2)} .
\label{eq:rate}
\end{equation}

A more robust choice is a least-squares fit of $\log e$ against $\log h$ (which is
how \code{estimate\_convergence\_rate(h, errors)} is implemented); the option
\code{use\_last=k} restricts the fit to the last $k$ levels in order to approximate
the asymptotic order.

\section{Measured Verification Results}
\label{sec:verify}

\begin{table}[htbp]
\centering
\caption{Measured convergence orders (manufactured solution, successively refined meshes)}
\label{tab:verify}
\begin{tabular}{@{}llccc@{}}
\toprule
Element & Mesh & measured $L^2$ & measured $H^1$ & theory $(L^2/H^1)$ \\
\midrule
P1 & triangles & 1.98 & 0.99 & 2 / 1 \\
P2 & triangles & 3.00 & 1.99 & 3 / 2 \\
Q1 & quadrilaterals & 2.00 & 1.00 & 2 / 1 \\
Q2 & quadrilaterals & 3.00 & 2.00 & 3 / 2 \\
VEM $k=1$ & triangles & 1.98 & 0.99 & 2 / 1 \\
VEM $k=1$ & quadrilaterals & 1.99 & 1.00 & 2 / 1 \\
VEM $k=1$ & pentagons & 2.00 & 1.00 & 2 / 1 \\
VEM $k=1$ & Voronoi polygons & 1.97 & 1.01 & 2 / 1 \\
P1 & tetrahedra & 1.95 & 0.98 & 2 / 1 \\
Q1 & hexahedra & 2.01 & 1.00 & 2 / 1 \\
VEM $k=1$ & hexahedra & 1.97 & 0.97 & 2 / 1 \\
RT0$\times$P0 & triangles (pressure) & 1.00 & --- & 1 \\
RT0$\times$P0 & triangles (flux) & 1.00 & --- & 1 \\
\bottomrule
\end{tabular}
\end{table}

To reproduce the table:

\begin{lstlisting}[style=shell]
& "D:\anaconda\envs\fealpy\python.exe" MosicaFE\tests\run_checks.py
\end{lstlisting}

\begin{remark}
The same script also prints the Helmholtz table for the constant-wavenumber case
$k=1$ (Table~\ref{tab:helmholtz-verification}): the measured orders are 1.98/0.99
(P1), 3.01/1.99 (P2), 2.00/1.00 (Q1) and 1.98/0.99 (VEM), i.e.\ identical to the
Poisson orders of Table~\ref{tab:verify}. The output also contains the resolution
indicator $kh$ together with $\lambda_{1,h}$, the discrete first Dirichlet
eigenvalue, so that the table can be checked against the well-posedness condition
$k^2<\lambda_1$ of Chapter~\ref{sec:helmholtz}.
\end{remark}

% ===========================================================================
\chapter{Visualization}
\label{chap:visual}

\section{Overview of the Plotting Tools}

\begin{table}[htbp]
\centering
\caption{Plotting functions in \module{postprocess}}
\label{tab:plots}
\begin{tabular}{@{}L{5.6cm}L{9.8cm}@{}}
\toprule
Function & Purpose \\
\midrule
\code{plot\_solution(u, mesh, V, ...)} & plot the solution. In 1D a curve; in 2D filled contours or a 3D surface;
in 3D a slice $z=\text{const}$. Polygonal meshes are triangulated by a fan from the cell centroids \\
\code{plot\_mesh(mesh, ...)} & plot the mesh (in 3D a mid-plane slice) \\
\code{plot\_error(u, mesh, space, exact, ...)} & plot the pointwise error $|u_h-u|$ \\
\code{plot\_convergence(h, errors, ...)} & log--log convergence curves with reference slopes $O(h^k)$ overlaid \\
\code{plot\_pressure\_2d(mesh, P, ...)} & piecewise constant RT0 pressure \\
\code{plot\_flux\_2d(mesh, Q, ...)} & flux arrow plot \\
\code{compare\_solutions(entries, ...)} & side-by-side comparison of several solutions \\
\bottomrule
\end{tabular}
\end{table}

Common keyword arguments: \code{title=}, \code{filename=} (saves a PNG),
\code{show=} (opens a window), \code{dpi=}, and \code{ax=} (draw into an existing
axes object, e.g.\ to build a panel figure).

\section{Plotting Solutions}

\begin{lstlisting}[language=MosicaFE,caption={Plotting 2D and 3D solutions}]
from MosicaFE import Mesh, FESpace, PoissonProblem
from MosicaFE.postprocess import plot_solution, plot_mesh, plot_error

# --- 2D contour ---
mesh = Mesh.rectangle(nx=32, ny=32)
V = FESpace.lagrange(mesh, 2)
u = PoissonProblem(V, f=source, g=lambda x: 0.0).solve()
plot_solution(u, mesh, V, title="FEM P2", filename="sol2d.png")

# --- 2D surface ---
plot_solution(u, mesh, V, mode="surface", filename="sol_surface.png")

# --- polygon mesh + VEM solution ---
mp = Mesh.voronoi_polygon(n_points=200, seed=3)
Vv = FESpace.vem(mp)
uv = PoissonProblem(Vv, f=source, g=lambda x: 0.0).solve()
plot_solution(uv, mp, Vv, title="VEM on polygons",
              filename="sol_vem.png")
plot_error(uv, mp, Vv, exact, filename="err_vem.png")

# --- 3D slice ---
m3 = Mesh.box(nx=8, ny=8, nz=8)
V3 = FESpace.lagrange(m3, 1)
u3 = PoissonProblem(V3, f=source3d, g=lambda x: 0.0).solve()
plot_solution(u3, m3, V3, z_slice=0.5, filename="sol3d_slice.png")
plot_mesh(m3, filename="mesh3d.png")
\end{lstlisting}

\section{Plotting Convergence Curves}

\begin{lstlisting}[language=MosicaFE,caption={Convergence curves}]
from MosicaFE.postprocess import plot_convergence

plot_convergence(h, errors, filename="convergence.png",
                 reference_rates=(1, 2, 3))
\end{lstlisting}

\begin{caution}
When \code{matplotlib} is not installed, all plotting functions merely print a notice
and return \code{None} instead of raising an exception, so batch scripts may call
them freely.
\end{caution}

% ===========================================================================
\chapter{Complete Examples}
\label{chap:examples}

\section{Example 1: Solving Poisson in Two Lines}

\begin{lstlisting}[language=MosicaFE,caption={One-line solution}]
import numpy as np
from MosicaFE import Mesh, solve_poisson

source = lambda x: 2*np.pi**2 * np.sin(np.pi*x[0]) * np.sin(np.pi*x[1])

u = solve_poisson(Mesh.rectangle(nx=32, ny=32), element="P2", f=source)
print(u.size)          # 4225
\end{lstlisting}

\section{Example 2: Changing the Element (One Code, Three Methods)}

\begin{lstlisting}[language=MosicaFE,caption={FEM / VEM / RT0 share the same workflow},label={lst:three}]
import numpy as np
from MosicaFE import FESpace, Mesh, MixedPoissonProblem, PoissonProblem

def exact(x):
    return np.sin(np.pi*x[0]) * np.sin(np.pi*x[1])

def grad(x):
    return np.array([np.pi*np.cos(np.pi*x[0])*np.sin(np.pi*x[1]),
                     np.pi*np.sin(np.pi*x[0])*np.cos(np.pi*x[1])])

def source(x):
    return 2*np.pi**2 * exact(x)

def zero(x):
    return 0.0

mesh = Mesh.rectangle(nx=32, ny=32)

# (a) classic FEM: P2 Lagrange
V = FESpace.lagrange(mesh, degree=2)
p = PoissonProblem(V, f=source, g=zero)
u = p.solve()
print("FEM P2 :", p.compute_errors(u, exact, grad))

# (b) VEM: only the space changes
V = FESpace.vem(mesh)
p = PoissonProblem(V, f=source, g=zero)
u = p.solve()
print("VEM k=1:", p.compute_errors(u, exact, grad))

# (c) RT0 x P0 mixed: space + problem class
V = FESpace.rt0(mesh)
p = MixedPoissonProblem(V, f=source)
Q, P = p.solve()
print("RT0-P0 :", p.compute_errors((Q, P), exact, lambda x: -grad(x)))
\end{lstlisting}

Output (on a 32$\times$32 mesh):

\begin{lstlisting}[style=textout]
FEM P2 : {'L2': 6.87e-05, 'Linf': 2.41e-03, 'H1': 7.38e-03}
VEM k=1: {'L2': 1.35e-03, 'Linf': 3.18e-03, 'H1': 1.09e-01}
RT0-P0 : {'L2': 1.64e-02, 'Linf': 4.55e-02, 'flux_L2': 6.30e-02}
\end{lstlisting}

\section{Example 3: Virtual Elements on a Polygonal Mesh}

\begin{lstlisting}[language=MosicaFE,caption={Voronoi polygonal mesh + convergence study}]
import numpy as np
from MosicaFE import FESpace, Mesh, PoissonProblem, convergence_study
from MosicaFE.postprocess import estimate_convergence_rate, plot_solution

def mesh_factory(level):
    return Mesh.voronoi_polygon(n_points=50 * 4**level, seed=7)

def problem_factory(mesh):
    return PoissonProblem(FESpace.vem(mesh), f=source, g=lambda x: 0.0)

h, errors = convergence_study(mesh_factory, problem_factory,
                              exact, grad, refinement_levels=3,
                              verbose=True)
print("L2 rate:", estimate_convergence_rate(h, errors["L2"]))

mesh = Mesh.voronoi_polygon(n_points=200, seed=7)
V = FESpace.vem(mesh)
u = PoissonProblem(V, f=source, g=lambda x: 0.0).solve()
plot_solution(u, mesh, V, title="VEM on Voronoi mesh",
              filename="vem_voronoi.png")
\end{lstlisting}

\begin{caution}
For random Voronoi meshes the maximum cell diameter is a very unstable choice for
$h$; using the mean cell diameter as the abscissa is preferable, which is achieved
simply by passing a custom \code{mesh\_factory} to \code{convergence\_study} (the
example uses \code{mesh.get\_mesh\_size()}; replace it inside the callback for a more
robust measure).
\end{caution}

\section{Example 4: The Low-Frequency Helmholtz Equation}

The Helmholtz equation $-\Delta u-k^2u=f$ is a built-in problem
(Chapter~\ref{sec:helmholtz}), so solving it costs exactly as much code as solving
the Poisson equation:

\begin{lstlisting}[language=MosicaFE,caption={Constant wavenumber, three equivalent spellings},label={lst:helmholtz-example}]
import numpy as np
from MosicaFE import (FESpace, HelmholtzProblem, Mesh, convergence_study,
                      estimate_convergence_rate, run, solve_helmholtz)

k = 1.0                                        # constant wavenumber
exact = lambda x: np.sin(np.pi*x[0]) * np.sin(np.pi*x[1])
grad = lambda x: np.array([np.pi*np.cos(np.pi*x[0])*np.sin(np.pi*x[1]),
                           np.pi*np.sin(np.pi*x[0])*np.cos(np.pi*x[1])])
source = lambda x: (2*np.pi**2 - k**2) * exact(x)   # manufactured solution

mesh = Mesh.rectangle(nx=32, ny=32)

# (a) one line
u = solve_helmholtz(mesh, element="P2", f=source, k=k)

# (b) step by step: mesh / element / solver / visualization
V = FESpace.lagrange(mesh, degree=2)           # or FESpace.vem(mesh)
problem = HelmholtzProblem(V, f=source, g=lambda x: 0.0, k=k)
u = problem.solve(method="direct")
print(problem.compute_errors(u, exact, grad))

# (c) by registry name
u = run("helmholtz", mesh, element="VEM", f=source, k=k)
\end{lstlisting}

Output on the 32$\times$32 mesh (P2, $kh=0.044$, $k^2/\lambda_{1,h}=0.051$):

\begin{lstlisting}[style=textout]
{'L2': 6.88e-05, 'Linf': 2.41e-03, 'H1': 7.38e-03}
\end{lstlisting}

The error norms are the ones of the Poisson problem on the same mesh, because in
this frequency range the two discrete operators differ by a bounded perturbation of
the mass matrix. What changes is the resolution requirement, and the example below
makes it visible: the same code, the same tolerance, but the mesh is refined until
$kh$ is small enough for the asymptotic order to show up.

\begin{lstlisting}[language=MosicaFE,caption={Convergence of the Helmholtz solver}]
def mesh_factory(level):
    return Mesh.rectangle(nx=4 * 2**level, ny=4 * 2**level)

def problem_factory(mesh):
    return HelmholtzProblem(FESpace.lagrange(mesh, 2), f=source,
                            g=lambda x: 0.0, k=k)

h, errors = convergence_study(mesh_factory, problem_factory, exact, grad,
                              refinement_levels=4, solver="direct")
print("L2 rate:", estimate_convergence_rate(h, errors["L2"], use_last=3))
\end{lstlisting}

\begin{lstlisting}[style=textout]
  level 0: h=0.35355  L2=3.5219e-02  H1=4.4393e-01  Linf=1.4645e-01
  level 1: h=0.17678  L2=4.4391e-03  H1=1.1632e-01  Linf=3.8060e-02
  level 2: h=0.08839  L2=5.5175e-04  H1=2.9416e-02  Linf=9.6074e-03
  level 3: h=0.04419  L2=6.8753e-05  H1=7.3752e-03  Linf=2.4076e-03

L2 rate: 3.006
\end{lstlisting}

Two numbers should be watched in such a table: the observed order ($3.006$, i.e.\
the P2 order of \eqref{eq:apriori}) and the resolution $kh$, which here runs from
$0.25$ down to $0.031$. What happens when the mesh is too coarse for the wave is
discussed in Chapter~\ref{chap:physics}; the verification data for $k=1$ on four
element families are collected in Table~\ref{tab:helmholtz-verification}.

\section{Example 5: The Reaction-Diffusion Equation}

\begin{lstlisting}[language=MosicaFE,caption={The built-in reaction-diffusion problem}]
import numpy as np
from MosicaFE import FESpace, Mesh, ReactionDiffusionProblem

c = 10.0
exact = lambda x: np.sin(np.pi*x[0]) * np.sin(np.pi*x[1])
source = lambda x: (2*np.pi**2 + c) * exact(x)

mesh = Mesh.rectangle(nx=32, ny=32)
V = FESpace.lagrange(mesh, degree=2)
problem = ReactionDiffusionProblem(V, f=source, g=lambda x: 0.0,
                                   reaction_coeff=c)
u = problem.solve()
print(problem.compute_error(u, exact, "L2"))     # ~6.8e-05
\end{lstlisting}

% ===========================================================================
\chapter{Advanced Topics}
\label{chap:advanced}

\section{Accuracy Pitfalls and Lessons Learned}
\label{sec:pitfall-vem}

While the three programs were being merged into a single library, four classes of
problems surfaced in practice. None of them was a ``small numerical accuracy
issue'': each of them produced wrong results or destroyed the convergence order.
They are recorded here both to explain why \mosica{} is written the way it is and to
help readers inspect their own codes.

\subsection{Pitfall 1: Dirichlet Elimination Without the Right-Hand-Side Correction}

\textbf{Symptom}: everything is fine for homogeneous boundary conditions, but the
solution is wrong for inhomogeneous ones.

\textbf{Diagnosis}: run the linear patch test (equation~\eqref{eq:patch}). When the
exact solution belongs to the discrete space, the numerical solution must be exact.

\textbf{Conclusion}: see equation~\eqref{eq:bc-lift}; step (1) cannot be omitted.
\mosica{} turns this into a test in \code{tests/test\_core\_spaces.py}.

\subsection{Pitfall 2: Insufficient Quadrature Order for Tensor-Product Elements}

\textbf{Symptom}: the computed $L^2$ convergence order for 3D hexahedral Q1 elements
is $-0.11$.

\textbf{Diagnosis}: the integrand of the Q1 stiffness matrix is a quadratic
polynomial in each direction, and with only one Gauss point per direction the rule is
exact only up to degree 1; the under-integration makes the cell matrices
non-positive-definite and systematically underestimates the energy.

\textbf{Conclusion}: use \code{max(2, 2*p)} as the quadrature order for
tensor-product elements (at least $p$ points per direction); see the caution after
Table~\ref{tab:quad}.

\subsection{Pitfall 3: Normalization of the Constant Mode in the VEM Projection}

\textbf{Symptom}: the VEM convergence order is correct on triangular and square
meshes, but on Voronoi polygonal meshes the $L^2$ order drops from 2 to below 0.5 and
the error stalls around $10^{-2}$.

\textbf{Diagnosis}: insert a linear exact solution and check whether the numerical
solution is exact (patch test); it fails on irregular polygons.

\textbf{Conclusion}: the constant mode must be normalized according to
equation~\eqref{eq:vem-const} so that the projection reproduces linear polynomials
exactly; otherwise the stabilization penalizes linear functions and consistency is
lost.

\subsection{Pitfall 4: Geometric Weights in the Load Vector / Mass Matrix}

\textbf{Symptom}: even after fixing the projection, the $L^2$ error on irregular
polygonal meshes is still too large.

\textbf{Diagnosis}: taking $\int_K\varphi_i=|K|/n$ for every vertex amounts to
assuming that each vertex ``owns an equal share'' of the cell. That is correct for
triangles and regular quadrilaterals, but for polygons with very different edge
lengths the vertices on short edges are strongly overestimated.

\textbf{Conclusion}: switch to the fan-wise linear interpolation of
equation~\eqref{eq:vem-load}, where the geometric weights are determined by the shape
of the cell itself. After the fix, the $L^2$ convergence order on Voronoi meshes
returns to 1.97.

\subsection{Pitfall 5: Repeated Indices in \code{numpy} Assignment}

\textbf{Symptom}: \code{b[dofs] += local} disagrees with the sparse matrix assembly,
and the row sums of the mass matrix do not match the load vector.

\textbf{Diagnosis}: when \code{dofs} contains repeated indices (for instance when
the mesh has duplicate vertices), the buffered assignment of \code{numpy} keeps only
the last value, whereas COO assembly sums them.

\textbf{Conclusion}: always use \code{numpy.add.at}.

\section{Mesh Refinement Strategies}

\begin{itemize}[itemsep=0.2em]
  \item \textbf{Uniform refinement}: \code{mesh.refine\_uniform()}, multiplying the
        number of triangles by 4 and of hexahedra by 8;
  \item \textbf{Adaptive refinement (planned)}: based on the residual-type error
        indicator
        \begin{equation}
        \eta_K=h_K\|f+\Delta u_h\|_{L^2(K)}
        +h_K^{1/2}\|[\partial_n u_h]\|_{L^2(\partial K)}
        \label{eq:estimator}
        \end{equation}
        cells with $\eta_K>\theta\max_K\eta_K$ are marked for subdivision. The
        \code{Mesh} of \mosica{} supports ragged cells, so local refinement is
        feasible as far as the data structure is concerned; what is missing is the
        refinement implementation.
\end{itemize}

\section{Performance and Vectorization}

\begin{enumerate}[label=(\arabic*),itemsep=0.2em]
  \item assembly accumulates COO triplets and converts them to CSR once, at the end;
  \item basis-function evaluations on the reference element are cached per quadrature
        point (\code{LagrangeSpace.\_quadrature\_data}), avoiding repeated work in
        every cell;
  \item for several physical problems on the same mesh, the result of
        \code{stiffness\_matrix()} can be reused;
  \item for large problems, switching to an iterative solver is recommended;
  \item for further speedup the cell loop can be vectorized with \code{numba} or by
        ``grouping cells by type and applying batched tensor operations''.
\end{enumerate}

\section{Relation to Other Libraries}

\mosica{} deliberately stays ``small, readable and modifiable'':

\begin{itemize}[itemsep=0.2em]
  \item compared with FEniCS / Firedrake, \mosica{} is not a general-purpose finite
        element platform but a \textbf{readable implementation for teaching and
        research}; the price is that only structured mesh generation and a handful of
        element families are supported;
  \item compared with \code{fealpy}, \mosica{} focuses more narrowly on ``a
        side-by-side comparison of three discretization approaches: FEM, VEM and
        RT0'';
  \item compared with the three original scripts, \mosica{} turns them into one set
        of interfaces, unifies meshes and solvers, and adds verification.
\end{itemize}

% ===========================================================================
\chapter{Extending the Framework}
\label{chap:extend}

\section{Adding a New Element Type}

The steps are:

\begin{enumerate}[label=(\arabic*),itemsep=0.2em]
  \item create \code{MosicaFE/spaces/yourspace.py} and derive from
        \code{BaseFESpace};
  \item implement \code{\_build\_dof\_map()}: build the map from cell-local degrees
        of freedom to global ones, and fill in \code{n\_dofs}, \code{dof\_support} and
        \code{dof\_is\_interior};
  \item implement \code{stiffness\_matrix()}, \code{mass\_matrix()} and
        \code{load\_vector(f)};
  \item implement \code{interpolate(f)} and \code{evaluate(cell, u, points)};
  \item register the new family in \code{FESpace.FAMILIES} in
        \code{spaces/factory.py};
  \item add a patch test and a convergence-order test under \code{tests/}.
\end{enumerate}

\begin{lstlisting}[language=MosicaFE,caption={Skeleton of an explicit-basis space},label={lst:newspace}]
from MosicaFE.spaces.base import BaseFESpace

class MySpace(BaseFESpace):
    family = "myspace"
    has_explicit_basis = True

    def _build_dof_map(self):
        # fill self.dof_map, self.n_dofs, self.dof_support, self.dof_is_interior
        ...

    def reference_nodes(self, cell_idx):
        # reference coordinates of local DOFs
        ...

    def eval_basis(self, xi):
        ...

    def eval_basis_grad(self, xi):
        ...

    def stiffness_matrix(self):
        return self._assemble(self._local_stiffness, symmetric=True)

    def mass_matrix(self):
        return self._assemble(self._local_mass, symmetric=True)

    def load_vector(self, f):
        return self._assemble_vector(lambda c: self._local_load(c, f))

    def interpolate(self, f):
        return self._evaluate_function(f, self.dof_coords)
\end{lstlisting}

\begin{important}
If the basis functions of the new element are explicit (higher-degree Lagrange
elements, DG elements, or other tensor-product families), \code{spaces/lagrange.py}
is a good model: it already implements the generic assembly loop ``evaluate on the
reference element, then transform cell by cell'', so usually only \code{eval\_basis}
and the degree-of-freedom layout have to be supplied.
\end{important}

\section{Adding a New Equation}

\begin{lstlisting}[language=MosicaFE,caption={Template for a new PDE},label={lst:newpde}]
from MosicaFE.physics.base import BasePhysics
from MosicaFE.physics.registry import register_problem

@register_problem("my_pde")
class MyProblem(BasePhysics):
    """Weak form: a(u, v) = l(v)."""

    def __init__(self, space, f=None, g=None, **kwargs):
        super().__init__(space, f=f, g=g, **kwargs)
        # store problem specific data here

    def assemble(self):
        K = self.space.stiffness_matrix()
        M = self.space.mass_matrix()
        A = K + 3.0 * M                      # example combination
        b = self.space.load_vector(self.f)
        return self.apply_boundary_conditions(A, b)
\end{lstlisting}

After registration, \code{available\_problems()} contains \code{"my\_pde"}, and both
\code{create\_problem("my\_pde", V, ...)} and
\code{run("my\_pde", mesh, element="VEM", ...)} are immediately available.

The built-in equations are written exactly after this template, and they are the
cheapest way to calibrate how much code a new equation really needs:
\code{PoissonProblem} and \code{ReactionDiffusionProblem} are 8 and 14 lines long,
and the Helmholtz problem of Chapter~\ref{sec:helmholtz} differs from the Poisson
problem by a single term, $A=\kappa K-k^2M$.

\begin{caution}
The template above assumes that the right-hand side \emph{is} the load vector of a
function, \code{space.load\_vector(f)}. Not every physical source fits that shape: a
point source $\delta(x-x_0)$, a boundary term $\int_{\Gamma_N}hv\dif s$, or a
discontinuous coefficient must be assembled separately and added to \code{b} before
\code{apply\_boundary\_conditions} is called. A complete worked example is
\code{examples/custom\_pde.py}, which registers
\code{PointSourcePoissonProblem($-\Delta u=f+\delta(x-x_0)$)}: for a nodal space the
identity $\varphi_i(x_0)=\delta_{ij}$ makes the Dirac functional a single entry
$b_j\mathrel{+}=1$.
\end{caution}

\section[Adding a New Mixed Formulation (RT0\texorpdfstring{$\times$}{x}P0 Template)]{Adding a New Mixed Formulation (Using RT0\texorpdfstring{$\times$}{x}P0 as a Template)}

A mixed formulation (a saddle-point problem) differs from the ``elliptic operator''
pattern above, but the division of labour is the same: \textbf{the space computes the
cell matrix blocks, and the physics layer only assembles them}. Inside the library,
RT0$\times$P0 is implemented exactly after this template and can be copied:

\begin{enumerate}[label=(\arabic*),itemsep=0.2em]
  \item if a new facet degree-of-freedom layout is needed, obtain the facet topology
        from \code{core/facets.py} first (global facet indices, \code{cell\_facets},
        \code{signs});
  \item implement the ``block'' interfaces in the space: here
        \code{flux\_mass\_matrix()} ($(\Phi,\Phi)$) and \code{divergence\_matrix()}
        ($(P,\Phi)$), together with \code{pressure\_load\_vector(f)};
  \item assemble the blocks into a saddle-point matrix in the physics layer with
        \code{bmat} (from \code{MosicaFE.\_compat}):
        \code{bmat([[M, -B.T], [B, None]])};
  \item to solve, use \code{"direct"}, \code{"schur"} (static condensation) or
        \code{Solver.minres} (symmetric indefinite);
  \item put the error analysis in \code{postprocess/error.py}, following
        \code{mixed\_error\_norms} (one set of norms for the pressure and one for the
        flux).
\end{enumerate}

\begin{important}
The \code{physics} layer contains only 6 lines of assembly code (see \code{assemble}
in \code{physics/mixed\_poisson.py}), while the cell geometry and the numerical
integration stay entirely in \code{spaces/rt0.py}. Replacing the mixed element (for
instance by BDM1, or by the $Q_2\times Q_1$ pair for Stokes) therefore only requires
providing the new blocks; the physics layer and the postprocessing can be reused
unchanged.
\end{important}

\section{Adding a New Mesh Type}

\begin{enumerate}[label=(\arabic*),itemsep=0.2em]
  \item add a class method to \code{core/mesh.py} returning
        \code{Mesh(nodes, elements, element\_type, name=...)};
  \item 2D cells are automatically reordered counterclockwise according to their
        signed area, and 3D tetrahedra according to their volume;
  \item if a 3D cell is neither a tetrahedron nor a hexahedron, supply the face table
        of every cell (a list of local vertex indices) through \code{faces=...}.
\end{enumerate}

\begin{lstlisting}[language=MosicaFE,caption={Custom cell shapes}]
import numpy as np
from MosicaFE import Mesh

# a two-element hexagonal domain (arbitrary polygons are fine)
nodes = np.array([[0.0, 0.0], [1.0, 0.0], [1.5, 0.9],
                  [0.5, 1.5], [-0.5, 0.9], [0.5, 0.6]])
elements = [[0, 1, 2, 5], [0, 5, 3, 4],
            [1, 2, 3, 5], [0, 3, 4, 5]]

mesh = Mesh.from_arrays(nodes, elements, boundary_nodes=[0, 1, 2, 3, 4])
print(mesh.summary())
\end{lstlisting}

\section{Development Workflow and Testing Requirements}

The recommended workflow is:

\begin{enumerate}[label=(\arabic*),itemsep=0.2em]
  \item write a \textbf{failing} test first (patch test, convergence order,
        symmetry, ...);
  \item then implement the feature;
  \item run \code{tests/run\_tests.py} and \code{tests/run\_checks.py};
  \item update the progress and verification records under \code{progress/}.
\end{enumerate}

Every new feature should cover at least:

\begin{itemize}[itemsep=0.2em]
  \item cell level: Kronecker property of the basis functions, partition of unity,
        symmetry / positive definiteness of the cell matrices;
  \item space level: symmetry of the global matrices, the constant null space, and
        the consistency of the row sums of $M$ with $b(f\equiv1)$;
  \item problem level: the patch test (a linear exact solution must be reproduced
        exactly) and the convergence order;
  \item end to end: the one-line interfaces \code{solve\_poisson} / \code{run} must
        work.
\end{itemize}

% ===========================================================================
\chapter{Troubleshooting and FAQ}
\label{chap:faq}

\section{Common Issues}

\subsection{Import Errors}

\textbf{Symptom}: \code{ModuleNotFoundError: No module named 'MosicaFE'}.

\textbf{Cause}: \mosica{} lives in \code{D:\textbackslash mosica\textbackslash
MosicaFE}, so the \textbf{parent directory} \code{D:\textbackslash mosica} has to be
added to the module path.

\begin{lstlisting}[style=shell]
# option 1: install in editable mode
cd D:\mosica\MosicaFE && python -m pip install -e .

# option 2: extend PYTHONPATH
$env:PYTHONPATH = "D:\mosica"
\end{lstlisting}

\subsection{Wrong Convergence Order}

Check the following, in order:

\begin{enumerate}[label=(\arabic*),itemsep=0.15em]
  \item run the patch test first (equation~\eqref{eq:patch} in
        Chapter~\ref{chap:physics}). If it fails, the problem is certainly in the
        assembly or in the boundary handling;
  \item check whether the quadrature order is sufficient (tensor-product elements
        are especially prone to under-integration);
  \item check whether the definition of \(h\) is consistent with the mesh refinement;
  \item check whether the exact solution is smooth enough (manufactured solutions
        should preferably use trigonometric functions or polynomials);
  \item for VEM, check that the stabilization option is \code{"trace"}.
\end{enumerate}

\subsection{Iterative Solver Does Not Converge}

\begin{itemize}[itemsep=0.2em]
  \item verify first with a direct method that the system itself is solvable;
  \item CG requires a symmetric positive definite matrix: saddle-point problems must
        use MINRES or a direct method;
  \item increase \code{maxiter} and relax \code{tol};
  \item if the condition number is very large, consider scaling or preconditioning.
\end{itemize}

\subsection{Out of Memory}

\begin{itemize}[itemsep=0.2em]
  \item make sure sparse matrices are being used (\code{scipy} installed);
  \item three-dimensional problems grow quickly: $n^3$ nodes correspond to roughly
        $6n^3$ tetrahedra;
  \item switching to an iterative method avoids the fill-in of direct
        factorizations.
\end{itemize}

\section{Frequently Asked Questions}

\paragraph{Q: How do I choose the polynomial degree?}
A: In the vast majority of cases P2 / Q2 offers the best accuracy-to-cost ratio; start
with P1 for a quick check, and switch to quadratic elements once the solution is
smooth. Virtual elements are currently available only for $k=1$.

\paragraph{Q: When should I use an iterative solver?}
A: Below roughly $2\times10^4$ degrees of freedom a direct solver is the least
trouble; for larger problems switch to CG (symmetric positive definite) or GMRES
(general).

\paragraph{Q: What exactly is the difference between VEM and FEM?}
A: On triangular meshes the $k=1$ virtual element space coincides with the P1 finite
element space, and the two give identical numerical results (in this library the
errors agree bit by bit). The advantage of VEM appears on \textbf{general polygonal
meshes}: it allows a different number of facets per cell and unstructured polygonal
partitions, whereas classical Lagrange finite elements require a uniform cell shape.

\paragraph{Q: Why does the RT0 solution return two arrays?}
A: The mixed formulation solves for the flux $q_h$ and the pressure $p_h$
simultaneously. \code{Q[c,i]} is the outward normal flux on the $i$-th facet of cell
$c$ (numbered so that facet $i$ is the one ``not containing vertex $i$''), and
\code{P[c]} is the constant pressure on that cell.

\paragraph{Q: Does RT0 require installing or linking an external program?}
A: No. RT0$\times$P0 is a native method of \mosica{}: the explicit basis functions
$\varphi_i=(x-a_i)/(d|K|)$ and the exact mass matrix are in \code{spaces/rt0.py}, the
facet topology and orientation are in \code{core/facets.py}, the saddle-point
assembly is in \code{physics/mixed\_poisson.py}, and the error norms are in
\code{postprocess/error.py}. The dependencies of the library are still just numpy
(scipy / matplotlib optional), and after \code{import MosicaFE} the call
\code{FESpace.rt0(mesh)} works directly.

\paragraph{Q: Can time-dependent problems be solved?}
A: The current version handles steady elliptic problems only. For time-dependent
problems the method of lines is recommended: space discretization gives
$M\dot u+Au=b$, and the time direction is advanced with an explicit or implicit
scheme of your choice. \code{mass\_matrix()} is already provided.

\paragraph{Q: Can nonlinear problems be solved?}
A: The framework targets linear problems. For nonlinear problems, write a Newton
iteration in the outer loop and use \mosica{} to assemble and solve the Jacobian
system at every step.

\paragraph{Q: How do I export results to ParaView?}
A: Use \code{space.nodal\_values(u)} to obtain nodal values, combine them with
\code{mesh.nodes} and \code{mesh.elements}, and write a \code{.vtu} or \code{.vtk}
file yourself; alternatively add an export function to \code{postprocess}.

\appendix

% ===========================================================================
\chapter{API Reference}
\label{app:api}

\section{Top-Level Module \module{}}

\begingroup\footnotesize
\begin{longtable}{@{}L{5.0cm}L{9.4cm}@{}}
\toprule
Name & Description \\
\midrule
\endfirsthead
\toprule
Name & Description \\
\midrule
\endhead
\code{Mesh} & mesh object and factory methods \\
\code{FESpace} & space factory: \code{lagrange} / \code{vem} / \code{rt0} / \code{create} \\
\code{FacetTopology}, \code{build\_facet\_topology} & facet topology and orientation (for $H(\mathrm{div})$ spaces) \\
\code{QuadratureRule} & quadrature rules on reference elements \\
\code{cell\_quadrature(...)} / \code{simplex\_quadrature(...)} & quadrature points and weights on physical cells / arbitrary simplices \\
\code{BaseFESpace} & base class of the spaces \\
\code{LagrangeSpace} / \code{VEMSpace} / \code{RT0Space} & the three space families \\
\code{rt0\_flux\_basis}, \code{rt0\_mass\_matrices} & explicit RT0 basis functions and exact mass matrices \\
\code{BasePhysics} & base class of the physics problems \\
\code{PoissonProblem} & $-\nabla\cdot(\kappa\nabla u)=f$ \\
\code{MixedPoissonProblem} & RT0$\times$P0 mixed formulation \\
\code{ReactionDiffusionProblem} & $-\Delta u+cu=f$ \\
\code{HelmholtzProblem} & $-\nabla\cdot(\kappa\nabla u)-k^2u=f$ (constant $k$) \\
\code{register\_problem}, \code{create\_problem}, \code{available\_problems} & PDE registry \\
\code{Solver} & linear solvers (direct / cg / gmres / bicgstab / minres / auto) \\
\code{compute\_error} / \code{compute\_errors} & error norms \\
\code{mixed\_error\_norms} & pressure / flux error norms for RT0$\times$P0 \\
\code{convergence\_study} & convergence studies \\
\code{estimate\_convergence\_rate} & convergence-order estimation \\
\code{plot\_solution} / \code{plot\_mesh} / \code{plot\_error} & visualization \\
\code{plot\_convergence} / \code{plot\_pressure\_2d} / \code{plot\_flux\_2d} & visualization \\
\code{solve\_poisson} & one-line Poisson solver \\
\code{solve\_helmholtz} & one-line Helmholtz solver \\
\code{run} & solve any registered PDE by name \\
\bottomrule
\end{longtable}
\endgroup

\section{\code{Mesh}}

\begingroup\small
\begin{longtable}{@{}L{5.0cm}L{9.4cm}@{}}
\toprule
Method / attribute & Description \\
\midrule
\endfirsthead
\toprule
Method / attribute & Description \\
\midrule
\endhead
\code{Mesh.interval(a,b,nx)} & one-dimensional interval mesh \\
\code{Mesh.rectangle(...)} & two-dimensional rectangular mesh (\code{triangle}/\code{quad}) \\
\code{Mesh.box(...)} & three-dimensional box mesh (\code{tet}/\code{hex}) \\
\code{Mesh.pentagon(nx,ny)} & pentagon mesh \\
\code{Mesh.voronoi\_polygon(...)} & Voronoi polygonal mesh \\
\code{Mesh.from\_arrays(...)} & custom mesh \\
\code{mesh.refine\_uniform()} & uniform refinement \\
\code{mesh.get\_mesh\_size()} & maximum cell diameter $h$ \\
\code{mesh.get\_cell\_vertices(i)} & coordinates of the cell vertices \\
\code{mesh.get\_cell\_measure(i)} & cell area / volume \\
\code{mesh.get\_cell\_diameter(i)} & cell diameter \\
\code{mesh.get\_cell\_centroid(i)} & cell centroid \\
\code{mesh.boundary\_nodes} & boundary vertex indices \\
\code{mesh.boundary\_facets} / \code{mesh.facets} & facet topology \\
\code{mesh.facet\_topology} & \code{FacetTopology}: global facet indices, shared facets, outward normal signs \\
\code{mesh.summary()} & one-line summary \\
\bottomrule
\end{longtable}
\endgroup

\section{\code{FESpace} and the Space Methods}

\begingroup\small
\begin{longtable}{@{}L{5.0cm}L{9.4cm}@{}}
\toprule
Method & Description \\
\midrule
\endfirsthead
\toprule
Method & Description \\
\midrule
\endhead
\code{FESpace.lagrange(mesh, degree)} & continuous Lagrange (P1/P2/Q1/Q2) \\
\code{FESpace.vem(mesh, degree=1)} & virtual elements ($k=1$) \\
\code{FESpace.rt0(mesh)} & RT0$\times$P0 mixed space \\
\code{FESpace.create(mesh, family, degree)} & create by name \\
\code{FESpace.available()} & available space families \\
\code{V.n\_dofs} & number of degrees of freedom \\
\code{V.get\_cell\_dofs(i)} & cell-local $\to$ global degrees of freedom \\
\code{V.get\_dof\_coordinates(k)} & coordinates of a degree of freedom \\
\code{V.boundary\_dofs()} & Dirichlet degrees of freedom \\
\code{V.eval\_basis(xi)} / \code{V.eval\_basis\_grad(xi)} & basis functions on the reference element \\
\code{V.stiffness\_matrix()} & stiffness matrix \\
\code{V.mass\_matrix()} & mass matrix \\
\code{V.load\_vector(f)} & load vector \\
\code{V.interpolate(f)} & interpolation \\
\code{V.evaluate(c, u, pts)}, \code{V.evaluate\_gradient(...)} & reconstruction and evaluation inside a cell \\
\code{V.nodal\_values(u)} & nodal values (for visualization) \\
\midrule
\code{V.flux\_mass\_matrix()} & RT0: the $(\Phi,\Phi)$ block (exact closed-form mass matrix) \\
\code{V.divergence\_matrix()} & RT0: the $(P,\Phi)$ block (divergence, always $\pm1$) \\
\code{V.pressure\_load\_vector(f)} & RT0: right-hand side $\int_K f$ of the pressure equation \\
\code{V.interpolate\_pressure(f)} / \code{V.interpolate\_flux(q)} & RT0: interpolation of each of the two fields \\
\code{V.evaluate\_flux(c, Q, pts)} & RT0: reconstruction of $q_h$ inside a cell \\
\code{V.locate\_cell(x)} & RT0: locate the cell containing a point \\
\code{V.n\_flux} / \code{V.n\_pressure} / \code{V.topology} & RT0: sizes and facet topology \\
\bottomrule
\end{longtable}
\endgroup

\section{Physics Problems}

\begingroup\small
\begin{longtable}{@{}L{5.0cm}L{9.4cm}@{}}
\toprule
Method & Description \\
\midrule
\endfirsthead
\toprule
Method & Description \\
\midrule
\endhead
\code{PoissonProblem(V, f, g, kappa)} & constructor \\
\code{problem.assemble()} & returns $(A,b)$ (boundary conditions included) \\
\code{problem.solve(method)} & assemble and solve \\
\code{problem.compute\_error(u, exact, norm)} & a single error norm \\
\code{problem.compute\_errors(u, exact, grad)} & all error norms \\
\code{HelmholtzProblem(V, f, g, k, kappa)} & constructor, $k$ and $\kappa$ are real constants \\
\code{helmholtz.k}, \code{helmholtz.kappa} & the stored coefficients ($A=\kappa K-k^2M$) \\
\code{solve\_helmholtz(mesh, element, f, g, k)} & one-line Helmholtz solve (rejects \code{"RT0"}) \\
\code{MixedPoissonProblem(...)} & returns $(Q,P)$ \\
\code{mixed.error\_norms(P, Q, p\_exact, q\_exact)} & pressure / flux errors \\
\code{mixed.check\_conservation()} & cell-wise mass conservation error \\
\code{mixed.flux\_at\_points(points)} & flux reconstruction at given points (cell location is automatic) \\
\code{mixed.cell\_flux(cell)} & flux vector at the cell centroid (for arrow plots) \\
\code{mixed.solve(method)} & three ways to solve the saddle-point system: \code{direct} / \code{schur} / \code{minres} \\
\bottomrule
\end{longtable}
\endgroup

\section{Solvers and Postprocessing}

\begingroup\small
\begin{longtable}{@{}L{5.0cm}L{9.4cm}@{}}
\toprule
Method & Description \\
\midrule
\endfirsthead
\toprule
Method & Description \\
\midrule
\endhead
\code{Solver.direct(A,b)} & direct method \\
\code{Solver.conjugate\_gradient(...)} & CG \\
\code{Solver.gmres(A,b,tol)} & GMRES \\
\code{Solver.bicgstab(A,b,tol)} & BiCGSTAB \\
\code{Solver.minres(A,b,tol)} & MINRES \\
\code{Solver.solve(A,b,method)} & unified entry point \\
\code{Solver.apply\_dirichlet(...)} & impose Dirichlet data \\
\code{Solver.recommend(n, symmetric)} & recommend a solver \\
\code{convergence\_study(...)} & convergence study \\
\code{estimate\_convergence\_rate(h, e)} & convergence order \\
\code{plot\_solution(u, mesh, V, ...)} & plot the solution \\
\code{plot\_mesh(mesh, ...)} & plot the mesh \\
\code{plot\_error(u, mesh, space, exact)} & plot the error \\
\code{plot\_convergence(h, errors)} & plot convergence curves \\
\code{plot\_pressure\_2d(mesh, P)} & plot the RT0 pressure \\
\code{plot\_flux\_2d(space=V, Q=Q)} & plot the flux (reconstructed from $q_h$ at cell centroids) \\
\code{compare\_solutions(entries)} & side-by-side comparison \\
\bottomrule
\end{longtable}
\endgroup

% ===========================================================================
\chapter{Mathematical Notation}
\label{app:notation}

\section{General Notation}

\begin{longtable}{@{}ll@{}}
\toprule
Symbol & Meaning \\
\midrule
\endfirsthead
\toprule
Symbol & Meaning \\
\midrule
\endhead
$\Omega$ & computational domain \\
$\partial\Omega$ & boundary of the domain \\
$d$ & spatial dimension (1, 2, 3) \\
$h$ & mesh size (maximum cell diameter) \\
$K$ & mesh cell \\
$\hat K$ & reference cell \\
$\kappa$ & diffusion coefficient (constants $\kappa>0$ in the built-in problems) \\
$k$ & wavenumber of the Helmholtz problem (constant; $k^2$ is the shift against $\lambda_1$) \\
$\lambda_1$, $\lambda_{1,h}$ & first Dirichlet eigenvalue of $-\Delta$, exact and discrete \\
$\lambda$ & generic eigenvalue (generalized problem $Kv=\lambda Mv$) \\
$n$ & number of degrees of freedom \\
$\mathcal{D}$ & set of Dirichlet degrees of freedom \\
\bottomrule
\end{longtable}

\section{Function Spaces}

\begin{longtable}{@{}ll@{}}
\toprule
Symbol & Meaning \\
\midrule
\endfirsthead
\toprule
Symbol & Meaning \\
\midrule
\endhead
$V$ & infinite-dimensional trial / test space \\
$V_h\subset V$ & finite-dimensional discrete space (finite or virtual elements) \\
$H^1(\Omega)$ & first-order Sobolev space \\
$H(\mathrm{div};\Omega)$ & space of vector fields with square-integrable divergence \\
$L^2(\Omega)$ & space of square-integrable functions \\
$\mathcal{P}_k$ & polynomials of degree at most $k$ \\
$Q_k$ & tensor-product polynomials of degree at most $k$ in each variable \\
\bottomrule
\end{longtable}

\section{Operators and Norms}

\begin{longtable}{@{}ll@{}}
\toprule
Symbol & Meaning \\
\midrule
\endfirsthead
\toprule
Symbol & Meaning \\
\midrule
\endhead
$\nabla u$ & gradient \\
$\nabla\cdot q$ & divergence \\
$\Delta u$ & Laplace operator \\
$\varphi_i$ / $\hat\varphi_i$ & basis functions on the physical / reference cell \\
$\Pi^{\nabla}$ & gradient projection operator of VEM \\
$\Pi_{\rm dof}$ & matrix representation of the projection on the degrees of freedom \\
$A_{\rm proj}$, $S_{\rm stab}$ & consistency and stabilization parts of the VEM cell matrix \\
$J_K$ & Jacobian matrix of the reference-to-physical mapping \\
$\|e\|_{L^2}$, $|e|_{H^1}$, $\|e\|_{L^\infty}$ & error norms \\
$q$, $p$ & flux and pressure in the RT0 mixed formulation \\
$A$, $b$ & matrix and right-hand side of the discrete linear system \\
\bottomrule
\end{longtable}

% ===========================================================================
\chapter{Bibliography}
\label{app:bib}

\begin{thebibliography}{9}
\bibitem{brenner} S. C. Brenner and L. R. Scott,
  \emph{The Mathematical Theory of Finite Element Methods}, 3rd ed.,
  Springer, 2008.
\bibitem{ern} A. Ern and J.-L. Guermond,
  \emph{Theory and Practice of Finite Elements}, Springer, 2004.
\bibitem{logg} A. Logg, K.-A. Mardal, and G. Wells (eds.),
  \emph{Automated Solution of Differential Equations by the Finite Element Method},
  Springer, 2012.
\bibitem{hitchhiker} L. Beir\~ao da Veiga, F. Brezzi, A. Cangiani,
  G. Manzini, L. D. Marini, and A. Russo,
  ``Basic principles of virtual element methods'',
  \emph{Mathematical Models and Methods in Applied Sciences}, 23(1):199--214, 2013.
\bibitem{vem2014} L. Beir\~ao da Veiga, F. Brezzi, L. D. Marini, and A. Russo,
  ``The Hitchhiker's Guide to the Virtual Element Method'',
  \emph{Mathematical Models and Methods in Applied Sciences}, 24(8):1541--1573, 2014.
\bibitem{vemstab} L. Beir\~ao da Veiga, F. Brezzi, L. D. Marini, and A. Russo,
  ``Virtual element methods for general second order elliptic problems on
  polygonal meshes'', \emph{SIAM Journal on Numerical Analysis}, 2013.
\bibitem{brezzi} F. Brezzi and M. Fortin,
  \emph{Mixed and Hybrid Finite Element Methods}, Springer, 1991.
\bibitem{rt} P.-A. Raviart and J.-M. Thomas,
  ``A mixed finite element method for 2nd order elliptic problems'',
  in \emph{Mathematical Aspects of Finite Element Methods}, Springer, 1977.
\bibitem{saad} Y. Saad, \emph{Iterative Methods for Sparse Linear Systems},
  2nd ed., SIAM, 2003.
\end{thebibliography}

\end{document}
